\documentclass[11pt]{article}
\usepackage{iftex}
\ifPDFTeX
  \usepackage[T1]{fontenc}
  \usepackage{lmodern}
\fi
\usepackage{amsmath,amssymb,amsthm}
\usepackage[margin=1in]{geometry}
\usepackage{booktabs,longtable,array}
\usepackage[numbers,sort&compress]{natbib}
\usepackage{microtype}
\usepackage{xurl}
\usepackage{hyperref}
\hypersetup{colorlinks=true,linkcolor=blue,citecolor=blue,urlcolor=blue,
  pdftitle={Goldbach-type representations with prescribed Liouville signs},
  pdfauthor={Jizhou Guo}}
\Urlmuskip=0mu plus 1mu
\setlength{\emergencystretch}{2em}
\newtheorem{theorem}{Theorem}[section]
\newtheorem{lemma}[theorem]{Lemma}
\newtheorem{proposition}[theorem]{Proposition}
\newtheorem{corollary}[theorem]{Corollary}
\newtheorem{inputthm}[theorem]{Input}
\theoremstyle{definition}
\newtheorem{definition}[theorem]{Definition}
\newcommand{\e}{\mathrm e}
\newcommand{\E}{\mathbb E}
\newcommand{\Prob}{\mathbb P}
\newcommand{\Z}{\mathbb Z}
\newcommand{\N}{\mathbb N}
\newcommand{\C}{\mathbb C}
\newcommand{\R}{\mathbb R}
\newcommand{\ind}{\mathbf 1}
\newcommand{\Dpret}{\mathbb D}
\newcommand{\ph}{\varphi}
\newcommand{\leanname}[1]{\nolinkurl{#1}}
\newcommand{\norm}[1]{\left\lVert#1\right\rVert}
\newcommand{\abs}[1]{\left\lvert#1\right\rvert}
\DeclareMathOperator{\rad}{rad}
\DeclareMathOperator{\tr}{tr}
\DeclareMathOperator{\HS}{HS}
\DeclareMathOperator{\Rea}{Re}
\DeclareMathOperator{\cond}{cond}
\title{Goldbach-type representations with prescribed Liouville signs}
\author{Jizhou Guo\\
  \href{mailto:mitsuha2021b@gmail.com}{\texttt{mitsuha2021b@gmail.com}}}
\date{October 2026}

\begin{document}
\maketitle
\begin{abstract}
Let $\lambda$ be the Liouville function. For every sufficiently large
integer $N$, each of the four ordered sign patterns
$(\lambda(a),\lambda(b))$, with $a,b\ge1$ and $a+b=N$, occurs
$N/4+O(N(\log N)^{-c})$ times, where one may take $c=10^{-200}$.
Equivalently, $\sum_{1\le n<N}\lambda(n)\lambda(N-n)
=O(N(\log N)^{-c})$; this proves, for the Liouville function,
the conjecture of Corradi and Katai that the sum is $o(N)$.
The proof combines the finite graph and residue-comparison estimates
of OpenAI's Lean~4 development of two-point correlations, the
Dirichlet zero-free half-plane $\Rea s>7/8$ from family~003 of the
same library, and the short-progression
variance bounds of Klurman--Mangerel--Ter\"av\"ainen; the deduction from
these inputs is carried out in the paper. A quotient-coordinate
construction fixes the graph multiplier while the total $N$ remains
in the progression modulus. Inserting the same half-plane into published
results of Pintz, Chen--Gupta--Li, Montgomery--Vaughan and Martin gives
further consequences: the exponent $3/5$ for the exceptional set in the
binary Goldbach problem, which Pintz proved under such a zero-free
hypothesis; the bound $O_\varepsilon(q^{7/3+\varepsilon})$ for the least
prime in a reduced class modulo $q$; the exponent $8$ in the
polylogarithmic bound for the least quadratic non-residue; and the bound
$O((\log p)^{24})$ for the least prime primitive root modulo $p$.
\end{abstract}

\section{Statement and context}
All logarithms are natural. Write $\lambda(n)=(-1)^{\Omega(n)}$ for $n\ge1$,
where prime factors are counted with multiplicity, and put
\begin{align*}
 C(N)&=\sum_{1\le n<N}\lambda(n)\lambda(N-n),\\
 A_{\epsilon_1,\epsilon_2}(N)&=
 \#\{(a,b)\in\N_{>0}^2:a+b=N,\ \lambda(a)=\epsilon_1,
                                      \lambda(b)=\epsilon_2\}.
\end{align*}
The pairs in this definition are ordered.

\begin{theorem}\label{thm:main}
There are absolute constants $C_0,N_0>0$ such that, for every integer
$N\ge N_0$ and every $(\epsilon_1,\epsilon_2)\in\{-1,1\}^2$,
\begin{equation}\label{eq:main}
 \abs{A_{\epsilon_1,\epsilon_2}(N)-\frac N4}
 \le C_0\frac{N}{(\log N)^{10^{-200}}}.
\end{equation}
Moreover,
\begin{equation}\label{eq:correlation}
 \abs{C(N)}\le C_0\frac{N}{(\log N)^{10^{-200}}}.
\end{equation}
\end{theorem}

The inputs are the named theorems of families 003 and 007 in
\cite{OpenAIMath}, at commit
\begin{center}\small
\nolinkurl{adc7f1241b42e322a6451854ab7e4b4c146bf78a}.
\end{center}
Family 003 supplies the Dirichlet zero-free half-plane; family 007 supplies
the finite closed-word estimates, residue comparisons and spectral transfer.
At this commit these Lean~4 theorems compile with only the axioms
\path{propext}, \path{Classical.choice} and \path{Quot.sound}.
Inputs~\ref{input:zf}, \ref{input:trace}--\ref{input:bad-spectral} and
\ref{input:paddingcost} transcribe the named declarations into mathematical
notation. The numerical exponent of Lemma~\ref{lem:explicitA}, the reductions
to these inputs and the proof of Theorem~\ref{thm:main} are written mathematics.
In particular, the Lean comparison theorems assert only
$\exists A_f\ge1000$; the bound $A_f\le10^6$ is read off from their
source proofs, and the value $c=10^{-200}$ depends on it.
The ancillary file \texttt{ExplicitBudget.lean} checks the polynomial
exponents $8457$ and $8458$ of the degree function. The other Lean files
cited below are also ancillary files of this submission.

The final two-point theorems of \cite{OpenAI007} fix affine forms
with positive leading coefficients and nonnegative constant terms
as the summation length grows, with constants depending on the forms.
Here the form $N-n$ varies with the summation length.
The argument runs the internal graph estimates at shift/length $1$,
with \cite{KMT} and Input~\ref{input:zf} supplying the variance of
$\lambda$ in progressions modulo $N$ needed for centring.

Theorem~\ref{thm:main} proves, for the Liouville function, the conjecture
of Corradi and Katai that $C(N)=o(N)$, as stated in
\cite[\S1.1]{Mangerel2024} and \cite[Conjecture~1]{Krishnamoorthy2026}.
Those papers give \cite{CK69} as its source.
De Koninck, Germ\'an and K\'atai \cite{DGK15} established a conditional
result under the existence of infinitely many Siegel zeros; we use the
description in \cite[\S1.1]{Mangerel2024}.
Mangerel \cite[Theorem~1.2]{Mangerel2024} proved $|C(N)|<N-1$ for every
$N\ge11$; for $N\ge2$, equality holds exactly when
$N\in\{2,3,5,10\}$.

For sign patterns, Shusterman asked whether every even $N>2$ is a sum
$a+b$ with $\lambda(a)=\lambda(b)=-1$; the question was posted by Pablo
on MathOverflow in August 2018 \cite{MO}.
Under GRH, Mangerel \cite[Theorems~1.2--1.3 and Proposition~1.4]{MangerelGRH}
proved this for all large even $N$, and all four patterns for every $N$
with a prime factor exceeding a fixed $p_0$. For sufficiently large
prime $N$, each pattern occurs $\gg N\exp(-C(\log\log N)^6)$ times.
His Remarks~1 and~9 replace GRH, for the two existence theorems, by
a zero-free rectangle
$\Rea s>1-(\log N)^{-\alpha}$, $|\Im s|\le(\log N)^3$,
for nonprincipal characters modulo $N$, with $0<\alpha<3/50$;
Input~\ref{input:zf} supplies this rectangle for large $N$.
CaptainSude \cite[Theorem~1.1]{CaptainSude} gave an elementary proof,
with a Lean formalization, for every even $N>2$.
Meng \cite[Theorem~A]{Meng2026} proved that all four patterns occur
for every $N\ge2$ outside $\{2,3,4,5,6,9,10\}$.
For $G_k(N)=\sum_{a_1+\cdots+a_k=N}\prod_i\lambda(a_i)$, with $a_i\ge1$,
Krishnamoorthy \cite[Theorems~1--2]{Krishnamoorthy2026} proved
$|G_k(N)|\ll_{B,k}N^{k-1}(\log N)^{-B}$ for each fixed $k\ge3$ and $B>0$,
and, for $k=2$,
$\#\{N\le X:|C(N)|>tN\}\ll_B t^{-2}X(\log X)^{-B}$ for $t>0$.
Theorem~\ref{thm:main} gives the asymptotic count of all four patterns
for every sufficiently large total. Constants and thresholds in this
note are chosen existentially; their numerical values are not tracked.

\begin{lemma}\label{lem:indicators}
For every $N\ge2$ and every pair of signs,
\begin{equation}\label{eq:indicator}
 A_{\epsilon_1,\epsilon_2}(N)
 =\frac{N-1+(\epsilon_1+\epsilon_2)
              \sum_{1\le n<N}\lambda(n)+\epsilon_1\epsilon_2 C(N)}4.
\end{equation}
Consequently a bound for $C(N)$ with an absolute logarithmic saving,
together with such a bound for $\sum_{n<N}\lambda(n)$, gives the four
counts with the same smaller saving. Conversely, the four counts imply
the corresponding bound for $C(N)$.
\end{lemma}
\begin{proof}
Expand
$\frac14\sum_{1\le n<N}(1+\epsilon_1\lambda(n))
                          (1+\epsilon_2\lambda(N-n))$.
The two linear sums agree under $n\mapsto N-n$.
For the converse use
$C(N)=\sum_{\epsilon_1,\epsilon_2}\epsilon_1\epsilon_2
 A_{\epsilon_1,\epsilon_2}(N)$; the four main terms cancel.
\end{proof}

\part{The reflected correlation}
\section{Zero-free input and progression estimates}

\begin{inputthm}[Family 003, {\cite[Theorem~1.1]{OpenAIQ78}}]\label{input:zf}
For every positive integer $q$, every Dirichlet character $\chi\pmod q$
and every $s\in\C$ with $\Rea s>7/8$, one has $L(s,\chi)\ne0$,
apart from the principal character's pole at $s=1$.
\end{inputthm}
The exact Lean name is
\begin{center}\small
\path{OAI.DirichletCharacter.LFunction_ne_zero_of_seven_eighths_lt_re}.
\end{center}
Its character is Mathlib's $\texttt{DirichletCharacter}\,\C\,q$, its modulus has
the instance \path{[NeZero q]}, and its explicit pole hypothesis is
$\neg(\chi=1\wedge s=1)$. The companion theorem
\path{OAI.riemannZeta_ne_zero_of_seven_eighths_lt_re} supplies the zeta
statement. Mathlib's zeta function is totalized at $1$; away from that
point the statement is the usual zero-free assertion for zeta.

Put $\e(t)=\exp(2\pi it)$ and, for 1-bounded multiplicative functions,
\[
 \Dpret(f,g;X)^2=\sum_{p\le X}\frac{1-\Rea(f(p)\overline{g(p)})}{p}.
\]

\begin{lemma}\label{lem:realmean}
There are absolute $C,X_0>0$ such that, for all $X\ge X_0$, all
$1\le q\le X$, with $q$ integral, and all real Dirichlet characters $\chi\pmod q$,
\begin{equation}\label{eq:realmean}
 \abs{\sum_{n\le X}\lambda(n)\chi(n)}
 \le C X(\log X)^{-1/40}.
\end{equation}
Principal and imprimitive characters are included.
\end{lemma}
\begin{proof}
For a nonprincipal primitive character, truncated Perron at a half-integer
$y\ge3$, with right line $1+1/\log y$ and height $T\asymp y^2$,
can be moved to $15/16$.
Montgomery--Vaughan's local expansion for $L'/L$
\cite[Lemma 11.1 and Theorem 10.17]{MV} and
Input~\ref{input:zf} give, uniformly on $15/16\le\sigma\le2$,
\[
 \abs{L'/L(\sigma+it,\chi)}\ll\log(2q(|t|+2)).
\]
Indeed every zero in the local expansion is at horizontal distance at
least $1/16$, and the number of local zeros is bounded by the stated
unit-height zero-counting theorem. The vertical integral is
$O(y^{15/16}\log(2qT)\log(2T))$; the horizontal integrals and Perron
error are smaller. Moving the cutoff by at most one von Mangoldt term
therefore gives
\begin{equation}\label{eq:psi}
 \psi(y,\chi)\ll y^{15/16}\log^2(2qy),\qquad y\ge2.
\end{equation}
An induced character differs from its primitive character by Euler-factor
terms of total size at most $O(\log(2q)\log(2y))$.
Removing prime powers costs $O(\sqrt y\log^2(2y))$.
Thus \eqref{eq:psi} also bounds
$\sum_{p\le y}\chi(p)\log p$, for every nonprincipal character, with an
absolute constant and without a lower bound on $y/q$.
The range $2\le y<3$ follows by the absolute bound for its finitely many
von Mangoldt terms.

Set $Y=(\log X)^{64}$. Partial summation, using $q\le X$, gives
\[
 \abs{\sum_{Y<p\le X}\frac{\chi(p)}p}
 \ll (\log X)^2Y^{-1/16}\ll1.
\]
For $f=\lambda\chi$, this and Mertens' theorem yield
\[
 \Dpret(f,1;X)^2
 =\sum_{p\le X}\frac1p+\sum_{p\le X}\frac{\chi(p)}p
 \ge\log\log X-O(\log\log\log X).
\]
For a principal character the second sum is nonnegative, so the same
lower bound holds.

Apply Matom\"aki--Radziwi\l\l--Tao
\cite[Lemma C.1, corrected version 3]{MRT} with its character slot equal
to the fixed modulus-one character and its real function slot equal to
the varying $f$. For $1\le|t|\le X$ it gives
$\Dpret(f,n^{it};X)\ge\frac14\sqrt{\log\log X}-O(1)$; for $|t|\le1$
it gives $\Dpret(f,n^{it};X)\ge\frac13\Dpret(f,1;X)-O(1)$.
Hence, after increasing an absolute $X_0$,
\[
 M:=\inf_{|t|\le\sqrt{\log X}}\Dpret(f,n^{it};X)^2
 \ge\frac1{32}\log\log X.
\]
The Hal\'asz inequality \cite[Lemma 7.4, modulus 1]{KMT}, with
$T=\sqrt{\log X}$, bounds $X^{-1}|\sum_{n\le X}f(n)|$ by
\[
 C\{(M+1)\exp(-M)+T^{-1/2}+(\log X)^{-1/4}\}
 \ll(\log X)^{-1/40}.
\]
The constants in the invocation of Lemma C.1 are absolute because its
character is fixed at modulus one.
\end{proof}

\begin{lemma}\label{lem:variance}
Fix $0<\nu\le1$. There are $C_\nu,X_\nu>0$ such that, for every
positive integer $q$, whenever $X\ge X_\nu$, $X\ge H\ge10q$, and
$R=H/q\ge\exp((\log X)^\nu)$,
\begin{equation}\label{eq:variance}
 \int_X^{2X}\sum_{a\bmod q}
   \abs{\sum_{x<n\le x+H\atop n\equiv a\ (q)}\lambda(n)}^2\,dx
 \le C_\nu\frac{XH^2}{q}(\log R)^{-1/1000}.
\end{equation}
The sum is over every residue class, including nonunits.
\end{lemma}
\begin{proof}
Write $\varepsilon=(\log R)^{-1/1000}$.
Use an external modulus upper bound which may exceed the actual $q$.
Set
\begin{gather*}
 T=\log X,\qquad \rho=\varepsilon^{11/10},\qquad \kappa=\rho/50,\\
 Q_*:=\left\lceil\max\{qR^\kappa,\ H\exp(-T^{2/5})\}\right\rceil,
 \qquad R_*=H/Q_*,\quad S_*=\log R_*.
\end{gather*}
For large $X$, uniformly in the stated range,
\begin{equation}\label{eq:kmtcap}
 \begin{gathered}
 q\le Q_*\le H/10,\\
 \tfrac12\min(R^{1-\kappa},\exp(T^{2/5}))
       \le R_*\le\min(R^{1-\kappa},\exp(T^{2/5})).
 \end{gathered}
\end{equation}
Both arguments of the maximum are $o(H)$ and the first is at least $q$.
The ceiling costs a factor at most two.
In particular, with $\nu_1=\min(\nu,2/5)$,
$S_*\ge T^{\nu_1}/3$ for large $X$ and $S_*\le T^{2/5}$.
The internal parameter lower bound is still
$\rho\ge S_*^{-1/200}$: before capping,
$S_*\asymp\log R$ and $0.0011<1/200$; after capping,
$\rho\ge T^{-0.0011}$ while $S_*^{-1/200}\asymp T^{-1/500}$.
Also $S_*\gg\rho^{-100}$ in both cases.

The explicit good-modulus set used in
\cite[Lemma 8.1]{KMT} is
\[
 \mathcal Q_{X,u,M}=
 \left\{q\le X:
 \begin{array}{l}
 L(s,\chi)\ne0\text{ for all }\chi\pmod q\text{ with }
                         \cond\chi>X^{u^{20}},\\
 \Rea s\ge1-M\log\log X/\log X,
                         \quad |\Im s|\le3X
 \end{array}\right\}.
\]
The deduction of Theorem 1.5 from Proposition 9.4 in version 5 uses
$\varepsilon^{1.1}$ in that proposition. Its good-modulus condition is
therefore $\mathcal Q_{X,\varepsilon^{6.6},\varepsilon^{-88}}$.
This is contained in
$\mathcal Q_{X,\varepsilon^6,\varepsilon^{-80}}$, which is contained in
$\mathcal Q_{X,\varepsilon^6,\varepsilon^{-36}}$ used for the real-function
reduction. The first set has the smaller conductor cutoff and larger
zero box. Since
\[
 \varepsilon^{-88}\frac{\log\log X}{\log X}
 \le (\log X)^{-0.912}\log\log X<\frac18,
\]
Input~\ref{input:zf} makes all three sets equal to all moduli up to $X$.
The reduction from Proposition 9.4 uses Lemma 7.7, which needs
$Z=\varepsilon^{-11}\le(\log X)^{1/20}$ and $q\le H/10$, both satisfied;
Lemma 6.2, the hybrid large sieve valid for all $q,T\ge1$; and
equation (50), which \cite{KMT} derive from Lemma~14 of Matom\"aki and Radziwi\l\l
\cite{MR2016}, which is reference [32] there.
We use this last step as printed. None restricts $q$ beyond
$q\le Q_*\le H/10$. The real-valued simplification is Corollary 1.6 and its proof.
Use $Q_*$ as the external $Q$ in these invocations. The accuracy
$\varepsilon$ and the actual modulus $q$ remain unchanged.

Here are the parameter margins for the nested invocations in that proof:
\begin{center}
\begin{tabular}{ll}
\toprule
Parameter & Margin\\
\midrule
$\varepsilon^{1.1}$ & $(\log R)^{-0.0011}\ge(\log R)^{-1/200}$\\
$\varepsilon^{6.6}$ & $\ge(\log X)^{-0.0066}>(\log X)^{-1/50}$\\
$20\varepsilon^{-220}\log\log X/\log X$ &
  $\le20(\log X)^{-0.78}\log\log X=o(1)$\\
$\varepsilon^{-11}$ & $\le(\log X)^{0.011}<(\log X)^{1/20}$\\
$\varepsilon^7$ & $\ge(\log X)^{-0.007}>(\log X)^{-1/13}$\\
\bottomrule
\end{tabular}
\end{center}
For the internal $\varepsilon^{1.1}$ invocation the last margin has
exponent $0.0077$, still below $1/13$.
We also check the auxiliary prime intervals in Section 9.3 of
\cite{KMT}, including its equation (52). Let $m\ge2$ be the smallest
integer with $m^{4m+2}S_*^m>T^{1/2}$. Define
\begin{align*}
 \log P_1&=\rho S_*,\qquad \log Q_1=S_*;\\
 \log P_j&=\rho j^{4j}S_*^j,\qquad
 \log Q_j=j^{4j+2}S_*^j\quad(2\le j<m),
\end{align*}
and $\log P_m=\rho^2T$, $\log Q_m=\rho T$.
For large $X$, $m\le b_\nu:=2+\lceil1/\nu_1\rceil$, since
$S_*\ge T^{\nu_1}/3$. With $\eta_a=1/100$, the required inequalities
are
\begin{equation}\label{eq:kmtintervals}
 \begin{aligned}
 \frac{\log\log Q_j}{\log P_{j-1}-1}&\le\frac{\eta_a}{4j^2},\\
 \frac{\eta_a}{j^2}\log P_j&\ge8\log Q_{j-1}+16\log j
 \quad(2\le j\le m).
 \end{aligned}
\end{equation}
For $j<m$, the first ratio is
$O_\nu(\log S_*/(\rho S_*))=o(1)$, using
$S_*\gg\rho^{-100}$; the second comparison gains a factor comparable
to $\rho S_*\to\infty$.
For $m=2$, its defining inequality gives $S_*>T^{1/4}/32$, so
$\log P_1\ge T^{1/4-11/10000}/32$.
Meanwhile $\log P_2\ge T^{1-22/10000}$ and
$\log Q_1\le T^{2/5}$; both terminal inequalities follow.
For $m\ge3$, minimality gives $\log Q_{m-1}\le T^{1/2}$.
The defining inequality for $m$ and its fixed bound $b_\nu$ give
$S_*^{m-1}\gg_\nu T^{(m-1)/(2m)}\ge T^{1/3}$.
Thus $\log P_{m-1}\gg_\nu T^{1/3-11/10000}$,
whereas $\log P_m\ge T^{1-22/10000}$ and
$\log\log Q_m\le\log T$. These bounds prove
\eqref{eq:kmtintervals} at the last interval as well.
For the first mean-value estimate in that proof, the case of $X_1$
on p.~36 of the arXiv version of \cite{KMT}, one needs $X/Q_1\ge q\mathcal T$, where
$\mathcal T=(X/H)R_*^{\rho/100}$ and $Q_1=H/Q_*=R_*$.
With $Q=q$, so that $R_*$ would equal $R$, this inequality fails by
the factor $R^{\rho/100}$;
the printed condition $Q_1\le h/(QP_1^{0.01})$ is inconsistent
with $Q_1=h/Q$. The buffer in $Q_*$ supplies the needed inequality:
\[
 q\mathcal T=\frac XR R_*^{\rho/100}
 \le\frac XR R^{\rho/100}
 \le\frac X{R_*}R^{-\rho/100}\le\frac X{Q_1}.
\]
The third inequality uses $Q_*\ge qR^{\rho/50}$.
The external $Q$ is a free parameter of Theorem~1.5 and Corollary~1.6,
whose right-hand sides involve $\varepsilon$, $\ph(q)$, $X$ and $H/q$;
thus this buffer preserves their conclusions.
The other uses of $Q_*$ also hold: $R_*^{-0.01\rho}\ll\varepsilon$
and $q/H=1/R\ll\varepsilon$ on p.~33 of the arXiv version, and $P_m\ge P_1$, since
$\log P_1\le T^{2/5}\ll\rho^2T=\log P_m$ on pp.~36--38 of that version.

The typicality in Definition 1.1 of \cite{KMT} means that
$\#\{p\mid q:p\le t\}\le\pi(t)/100$ for every $t\ge y$.
Take $y=R_*^{\varepsilon^2}$.
There are at most $\log X/\log2$ distinct prime factors of $q$, whereas
$\pi(y)/\log X\to\infty$: before capping,
$\log y\asymp(\log R)^{0.998}$, and after capping
$\log y\gg T^{0.398}$.
All larger cutoffs satisfy the same condition. The cutoff
$R_*^{\varepsilon^{1.1}}$ in Proposition 9.4 is larger than $y$.

Corollary 1.6, restricted to unit classes, now gives the variance about
its real-character main term with error
$O(\varepsilon\ph(q)X R^2)$. A nonreal main term may be omitted in
that corollary. For a real main term,
Lemma~\ref{lem:realmean} at cutoff $3X$ gives the relative squared cost
\[
 O\bigl((q/\ph(q))^2(\log X)^{-1/20}\bigr)=o(\varepsilon),
\]
using $q/\ph(q)\ll\log\log(3q)$.

To include nonunits, group the classes by $u=(a,q)$ and put
$n=um$, $x=ut$, $q'=q/u$, $X'=X/u$, $H'=H/u$.
Then $H'/q'=R$, $R\le X'\le X$, $q'\le X'$, and $dx=u\,dt$.
At every such scale use the same $\varepsilon=(\log R)^{-1/1000}$.
Choose its own $Q_*$ by the displayed formula with $X',H',q'$.
The zero-box and all lower-parameter checks are uniform for
$X'\in[R,X]$: their worst lower scale is $X'=R$, where, for example,
$\varepsilon^{-88}\log\log R/\log R=o(1)$.
The typicality argument is repeated at that scale using
$\omega(q')\le\log X'/\log2$ and the corresponding $y$; one still has
$\log R\ge(\log X')^\nu$.
The function $(\log\log t)^2(\log t)^{-1/20}$ is decreasing for large
$t$, so its value at $X'$ is at most its value at $R$, which is
$o(\varepsilon)$.
Complete multiplicativity and $|\lambda(u)|=1$ show that this gcd group
contributes at most $C_\nu\varepsilon\ph(q/u)X R^2$.
Finally $\sum_{u\mid q}\ph(q/u)=q$, which gives \eqref{eq:variance}.
\end{proof}

\begin{lemma}\label{lem:fourier}
Fix $0<\nu_0\le1$. For sufficiently large $N$ and every real $D$ with
$\exp((\log N)^{\nu_0})\le D\le N$, one has
\begin{equation}\label{eq:fourier}
 \sup_{\theta\in\R}\frac1N\sum_{r=1}^N
   \abs{\sum_{0\le k<\lfloor D\rfloor}\lambda(Nk+r)\e(k\theta)}
 \ll_{\nu_0}D(\log D)^{-1/3000}.
\end{equation}
\end{lemma}
\begin{proof}
Put $K_D=\lfloor D\rfloor$, $Q_0=\lfloor\sqrt D\rfloor$.
Dirichlet approximation gives $\theta=a/q_0+\beta$ with
$1\le q_0\le Q_0$ and $|\beta|\le1/(q_0Q_0)$.
Set
\[
 h=\left\lfloor\min\{D/(\log D)^{10},q_0Q_0/100\}\right\rfloor,
 \qquad t_0=h/(\log D)^{20}.
\]
For $t_0\le t\le h$, one has
$t/q_0\gg\sqrt D/(\log D)^{30}$, and $|\beta|h\le1/100$.
For $0\le v<K_D$ define
\[
 U_{r,v}(t)=\sum_{1\le j\le t}\lambda(N(v+j)+r)
                                     \e(a(v+j)/q_0).
\]
Group by $v+j\pmod{q_0}$ and apply Cauchy--Schwarz.
The map $(r,b)\mapsto Nb+r$ is a bijection onto all classes modulo
$Q=Nq_0$ for $1\le r\le N$, $0\le b<q_0$.
Moving the starting point $Nv+r$ to $Nv+u$, $0\le u<N$, changes each
progression sum by at most two terms, since $Q\ge N$.
Integrating in $u$ gives
\begin{equation}\label{eq:startaverage}
 \sum_{r,v}|U_{r,v}(t)|^2
 \ll\frac{q_0}{N}\int_0^{NK_D}\sum_{b\bmod Q}
   \abs{\sum_{x<n\le x+Nt\atop n\equiv b\ (Q)}\lambda(n)}^2dx
   +NK_Dq_0^2.
\end{equation}
For $0\le x<Nt$, each class has $O(t/q_0+1)$ terms, giving total cost
$O(Nt^3)$ after multiplication by $q_0/N$.
Cover the rest by dyadic intervals with starting scale $X'\ge Nt$;
their starting scales have sum $O(ND)$.
Lemma~\ref{lem:variance} applies with length $Nt$ and modulus $Q$.
One may take $\nu=\nu_0/2$, since for large $N$
\[
 (2\log N+1)^{\nu_0/2}
 \le \tfrac12(\log N)^{\nu_0}-30\log\log N-O(1)
 \le\log(t/q_0).
\]
Thus, with $\varepsilon_D=(\log D)^{-1/1000}$,
\begin{equation}\label{eq:windowvariance}
 \frac1{NK_D}\sum_{r,v}|U_{r,v}(t)|^2
 \ll_{\nu_0}\varepsilon_Dt^2+t^3/D+q_0^2
 \ll_{\nu_0}\varepsilon_Dt^2.
\end{equation}
The final inequality uses $t/D\le(\log D)^{-10}$ and
$q_0/t\ll(\log D)^{30}/\sqrt D$.

Abel summation inserts $\e(\beta j)$ into the window.
At each $t\ge t_0$, average first in $(r,v)$ and use
\eqref{eq:windowvariance}; the mean absolute value is
$O_{\nu_0}(\sqrt{\varepsilon_D}\,t)$.
For $t<t_0$ use $|U_{r,v}(t)|\le t$.
The mean absolute value of the full $\theta$ window is consequently
$O_{\nu_0}(h\sqrt{\varepsilon_D}+h(\log D)^{-40})$.

For any 1-bounded sequence $(z_k)$ the deterministic identity
\[
 \frac1h\sum_{v=0}^{K_D-1}\sum_{j=1}^h z_{v+j}
 =\sum_{k=0}^{K_D-1}z_k+O(h+1)
\]
follows by counting the multiplicity of each index: the sum of the
absolute differences of the coefficients is at most $h+1$.
Use it with $z_k=\lambda(Nk+r)\e(k\theta)$ and average in $r$.
The resulting bound is
$O_{\nu_0}(D\sqrt{\varepsilon_D}+h+1)$, which is
$O_{\nu_0}(D(\log D)^{-1/2000})$ and implies \eqref{eq:fourier}.
The frequency is fixed throughout each average, so the supremum remains
outside the $r$ average.
\end{proof}

\section{Finite graph inputs and their parameters}
Throughout this section the prefix of a family 007 Lean name is
\path{OAI.TwoPointCorrelations}. Thus the names printed below specify
fully qualified declarations in that namespace.

\begin{definition}[Arithmetic word data]\label{def:words}
A prime family consists of disjoint finite prime sets $\mathcal P,
\mathcal Q$ and a finite supply $\mathcal S$ of pairs $(d,q)$, where $d,q$
are squarefree, $d$ has $J$ prime factors in $\mathcal P$, and $q$ has at
most $m_Q$ prime factors in $\mathcal Q$. At multiplier $h\ge1$, every
$p\in\mathcal P$ is required not to divide $h$.
This is the mathematical data of \leanname{ProhibitedPrimeFamily}.
For a signed step $a=(\epsilon,d,q)$, $\epsilon\in\{-1,1\}$, put
$\Delta_h(a)=\epsilon hqd$ and
$\Delta_h(w)=\sum_{a\in w}\Delta_h(a)$ for a word $w$.
A word is positive at $x\in\Z$ when, for each successive step, its whole
divisor $qd$ divides the departure integer $x+\Delta_h(w_{<i})$.
\end{definition}

For completeness we specify the prohibited predicate used in the graph.
A word $w=(a_0,\ldots,a_{R-1})$ is forward-prohibited if
$3\le R\le s$, all its pairs belong to $\mathcal S$, consecutive tuple
labels differ, and each tuple prime occurs on an interval of indices.
There must also be a prime $p$ in the first tuple but not in the last,
and an index $0<a<R-1$, such that
$p\mid\Delta_h(w_{\ge a})$.
A forward-prohibited word is minimal if none of its shorter contiguous
subwords, in either orientation, is forward-prohibited. Reversal changes
the order of the steps and reverses every direction.
An integer is prohibited if some minimal forward-prohibited word is
positive there. This is exactly \leanname{ProhibitedSite}, built from
\leanname{PositiveWord}, \leanname{ForwardProhibited}, and
\leanname{MinimalWord}. Positivity here is a divisibility condition for
integers of either sign.

The residue model assigns independent uniform variables
$R_p\in\Z/p\Z$ for $p\in\mathcal P\cup\mathcal Q$.
Integer divisibility at $x+u$ is replaced by $R_p+x+u=0\pmod p$.
A common CRT lift of these variables defines every word predicate; the
choice of lift has no effect. All expectations in the remainder of this
section are with respect to this finite probability space.

\begin{inputthm}[Closed-word column bound]\label{input:trace}
Fix $h\ge1$, $C\ge0$ and $W\ge1$. For all sufficiently large $L$, the
following holds for every prime family of Definition~\ref{def:words}.
Suppose its center set is the union of pairwise disjoint bands
$\mathcal P_1,\ldots,\mathcal P_J$, with masses
$V_i=\sum_{p\in\mathcal P_i}p^{-1}$, and suppose
\begin{gather}
 J\ge1,\quad J+m_Q\le C\log L,\quad m_Q\le100\log L,
 \quad1\le V_i\le\min(2W,L^2),\label{eq:tracehyp1}\\
 1\le\sum_{p\in\mathcal P}p^{-1}\le L^2,
 \qquad\sum_{p\in\mathcal Q}p^{-1}\le L^2,\label{eq:tracehyp2}\\
 \exp(L^{199/200})\le H\le p\le Y\le\exp L
       \quad(p\in\mathcal P),\qquad
 p\le B\le\exp L\quad(p\in\mathcal P\cup\mathcal Q),\label{eq:tracehyp3}\\
 q\le Q_{\max}\le\exp(100L+1)\quad((d,q)\in\mathcal S),
 \qquad\prod_{i\ne j}p_i\le D_{\max}\le\exp(2L)
       \quad(p_i\in\mathcal P_i).\label{eq:tracehyp4}
\end{gather}
Set $k=\lfloor L\rfloor$, $s=\lfloor L^{1/10}\rfloor$, and
$D_b=2k+\lfloor L^{1/12}\rfloor s$.
Let $\mathcal D$ be the tuple products, and take any finite squarefree
padding-divisor set supported on $\mathcal Q$.
Use a nonnegative padding coefficient $u(q)\le4^{\omega(q)}$, a
positive Q-residue-dependent vertex weight $g$ with $g^2\ge1$, a row cap
$K\ge1$, and Q-residue-dependent extra cuts. Retain precisely the
eligible edges, the row-density cap, these extra cuts, and the
nonprohibited vertices. Center each tuple by
$\prod_{p\mid d}(\ind_{p\mid x}-1/p)$.

For any finite state set mapped into tuple labels times integer sites,
any deterministic edge gates, any fixed initial state and any prescribed
sequence of $2k$ directions, let $\mathcal F$ be the numerical closed-word
catalog obtained from pairs of $k$-step paths with distinct consecutive
tuple labels in each half. No extra condition is imposed at the seams.
If $Z_w$ is the actual signed masked scalar word weight, then
\begin{equation}\label{eq:wordmajorant}
 \exp(108L)\sum_{w\in\mathcal F}\abs{\E Z_w}
 \le\{K(2\exp(150)\sqrt W)^J\}^{2k}.
\end{equation}
The state-to-site map need not be injective in this column statement.
\end{inputthm}

The exact name is \leanname{eventually_actual_column_trace}.
The numerical catalog is its \leanname{actualClosedPairCatalog}, with
\leanname{closedTraceFiber} fixing the direction sequence.
The theorem takes the threshold in $L$ after $h,C,W$ are fixed and
before all prime families, numerical words, gates and state maps.
The more general theorem \leanname{eventually_prohibited_column_trace_total}
is the source of this column bound.
The correspondence between the signed masked word and the centered
residue expectation is
\leanname{maskedClosedWord_residue_average}: for a closed word of
positive length $R\le D_b$, the residue average of the actual masked
word equals the centered product times its attached-catalog avoidance
factor and retained departure weight. Its hypotheses are disjoint prime
bands, the matching label assignment and a nowhere-zero $g$.
On a closed walk the square-root vertex denominators combine into the
departure squares; the idempotent vertex masks combine into the attached
avoidance condition. This equality retains the signs before taking the
absolute values in \eqref{eq:wordmajorant}.

\begin{inputthm}[Finite comparisons]\label{input:comparison}
There is an absolute integer $A_f\ge1000$ such that, for all sufficiently
large $L$, the following comparisons are uniform after $L$ is fixed.
The coordinate moduli are positive and pairwise coprime, number at most
$\exp L+1$, and are at most $\exp L$.
A Boolean test can be any unary test of a specified residue coordinate.
A depth-at-most-20 circuit of size at most $\exp(L^6)$ has empirical
versus independent-residue error at most $\exp(-L^{10})$ on any integer
interval of length $U\ge\exp(L^{A_f}/2)$.

In the following actual-word and positive-cost comparisons we use
$u(q)=4^{\omega(q)}$ and $g(n)^2=5^{\omega_{\mathcal Q}(n)}$,
with the Q-degree cut $\omega_{\mathcal Q}(n)\le400\log L$.
These comparisons concern these arithmetic weights and the positive
deletion atoms defined in Section~4.
For an actual closed masked word the error is at most $\exp(-L^9)$.
Its word length $R$ satisfies $0<R$, $R+1\le4L$, $RJ\le L^2$;
$s\le L$, $J+m_Q\le L^2$, the supply has size at most $\exp(101L)$,
and every padding choice has at most $100\log L$ prime factors.
The Q-degree cut is at $\lfloor400\log L\rfloor$.
The comparison holds on $a+\ell\{0,\ldots,U-1\}$ for arbitrary
$a,\ell\ge0$ with $(\ell,p)=1$ for every pool prime.

Positive deletion atoms satisfy an error of at most
$2\exp(-L^9)$, uniformly in the prime family and any integer offset.
The average degree-truncated padding weight is at most its product-law
normalizer plus $\exp(-L^9)$, and at most twice that normalizer.
Sums of atoms incur the sum of the individual absolute errors.
Bin-dependent interval lengths may be used, provided each is at least
$\exp(L^{A_f}/2)$.
\end{inputthm}

The first statement is
\leanname{BravermanDepth22Input.eventually_residue_comparison}, obtained
from \leanname{BravermanDepth22Input.finite_residue_comparison}.
The actual-word statement is
\leanname{BravermanDepth22Input.eventually_actual_affine_word_comparison}.
The last statements are the specializations of
\begin{quote}\small
\path{BravermanDepth22Input.eventually_variable_padding_comparison},\\
\path{BravermanDepth22Input.eventually_prohibited_row_comparison},\\
\path{BravermanDepth22Input.eventually_variable_tuple_degree_deletion},\\
\path{BravermanDepth22Input.eventually_actual_padding_weight}.
\end{quote}
The first, unweighted comparison is stated for consecutive origins.
Its arbitrary unary tests can be precomposed with
$z\mapsto a+\ell z$; multiplication by an invertible $\ell$ preserves
the independent uniform law. The named affine word theorem performs
this precomposition for the full weighted word.
There is no upper bound on $a,\ell$, nor a requirement that the product
of all pool primes be small. The proved premise
\leanname{bravermanDepth22Input} supplies the circuit comparison.

The weighted comparison expands at most $4L$ active-state slots, each
with at most $\exp L$ coordinates and at most $400\log L$ active bits,
and at most $L^2$ center bits. The coefficient table is bounded by
$\exp(L^4)$, and the surviving-site circuit has depth at most 19 and
size at most $\exp(L^3)$. This is the exact budget in
\leanname{BravermanDepth22Input.eventually_affine_weighted_word_comparison}.
It gives the actual-word and positive-atom versions above. Arbitrary
offsets change the unary residue tests and preserve these budgets.

\begin{inputthm}[Prohibited sites and spectral transfer]\label{input:bad-spectral}
For every fixed $C\ge0$, for all sufficiently large $L$, uniformly in
the multiplier $h$, the prime family and its supply,
\begin{equation}\label{eq:badprob}
 \Prob(\text{a residue-model site is prohibited})
 \le\exp(-\tfrac12 L^{199/200})
\end{equation}
if $1\le s\le L^{1/10}$, $J+m_Q\le C\log L$, the pool masses are
at most $L^2$, the center mass is at least 1, and all center primes
are at least $\exp(L^{199/200})$, with all pool primes at most $\exp L$.

For any finite-dimensional complex Hilbert space, finite self-adjoint
family $(B_d)$, orthogonal projection $E$, and nonnegative real
numbers $a,b$, suppose
\[
 \norm{B_d}\le a,\qquad
 \sum_d\norm{B_dv}^2\le b^2\norm v^2\quad(Ev=v).
\]
If $\mathcal H$ is its nonbacktracking operator, then
\begin{equation}\label{eq:spectralinput}
 \norm{E\sum_d B_d E}\le3\max\{a,b,\rad(\mathcal H)\}.
\end{equation}
\end{inputthm}
The exact names are \leanname{eventually_prohibited_density} and
\leanname{noncommuting_spectral_transfer_max}; the common-bound version
is \leanname{noncommuting_spectral_transfer}.
The probability bound follows from the finite theorem
\leanname{ProhibitedPrimeFamily.probability_bound}, which includes the
whole prohibited catalog. This probability bound concerns a single site. In Lemma~\ref{lem:rare}
it is combined by Cauchy--Schwarz with second moments of the positive
weights, and sums over many sites are compared with the residue model
by Input~\ref{input:comparison}.
The nonbacktracking operator uses distinct successive tuple labels;
different $B_d$ need not commute, and $E$ need not commute with them.

\begin{lemma}[A numerical finite-comparison exponent]\label{lem:explicitA}
Every finite comparison used here may be run with $A_f=1000000$.
\end{lemma}
\begin{proof}
We record the exponent extraction from the explicit circuit budgets in
\cite{OpenAIMath}. Write $\mathcal G(f,c)$ for
$\exists K\ \forall j\ge0:\ f(j)\le K(j+1)^c$.
The proof of \leanname{bravermanDegree_polynomial} uses the rules
\[
 \mathcal G(f,c),\mathcal G(g,d)
 \Longrightarrow\mathcal G(f+g,c+d),\quad
 \mathcal G(fg,c+d),\quad\mathcal G(f^r,rc).
\]
Constants have exponent 0 and $j$ has exponent 1.
Applying these rules in the order of the source definitions gives
\begin{center}
\begin{tabular}{@{}p{.56\textwidth}p{.40\textwidth}@{}}
\toprule
Budget & Polynomial exponent in $j+1$\\
\midrule
\leanname{bravermanBase} & 2\\
\leanname{bravermanNormExponent} & 45\\
\leanname{bravermanErrorExponent} & 1\\
\leanname{bravermanSwitches} & 47\\
\leanname{switchingDenominator} at that switch count & 94\\
\leanname{switchingDegree} at depth 89 & $47+89\cdot94=8413$\\
\leanname{bravermanDegree} & $44+8413=8457$\\
\bottomrule
\end{tabular}
\end{center}
Increasing the positive-growth exponent by one, as in
\leanname{bravermanDegree_polynomial_positive}, gives a fixed bound
$K(j+1)^{8458}$ for some absolute $K>0$.
The theorem
\path{ReflectedLiouville.braverman_degree_positive_explicit_exponent}
in the ancillary file \texttt{ExplicitBudget.lean} proves this numerical exponent
for the source's literal degree function.
The dyadic comparison
\leanname{AC0Circuit.dyadic_comparison} requires independence degree
at least this explicit \leanname{bravermanDegree}.
The logarithmic-budget conversion changes $K$, not 8458.

In the finite-residue comparison take bit depth
$b_L=\lceil L^{15}\rceil$, accuracy $\varepsilon_L=\exp(-L^{11})$,
and independence order
$t_L=\lceil L^{17\cdot8458+1}\rceil=\lceil L^{143787}\rceil$.
The encoded size is at most $\exp(L^{17})$; hence, for large $L$,
its required independence order is at most $t_L$.
The total number of encoded bits is at most $\exp(2L)$.
The low-order CRT/Fourier error is bounded by
\[
 \frac{(\exp L)^{t_L}}{U}(t_L+1)(\exp(2L))^{t_L}
 \le\exp(-L^{11})
 \quad\text{if }U\ge\exp(L^{1000000}/2).
\]
The two decoder errors total at most
$2\exp(3L)2^{-b_L}$, so the finite formula gives error at most
$\exp(-L^{10})$.
The signed-state expansions and all positive costs have total
coefficient mass at most $\exp(C L^5)$ for a fixed $C$; this changes the
error to at most $\exp(-L^9)$ and does not change the required interval
exponent. In the source,
\leanname{BravermanDepth22Input.eventually_residue_comparison} takes
$T:=17C+1$ and $A:=T+1002$, where $C$ is the exponent of
\leanname{bravermanDepth22Input}. Here $C=8458$, so
$A=17\cdot8458+1003=144789$.
The comparisons
\leanname{eventually_scalar_residue_comparison},
\leanname{eventually_active_state_comparison},
\leanname{eventually_weighted_word_comparison},
\leanname{eventually_affine_weighted_word_comparison},
\leanname{eventually_actual_affine_word_comparison},
\leanname{eventually_positive_weight_comparison},
\leanname{eventually_positive_prime_cost_comparison},
\leanname{eventually_prohibited_cost_comparison},
\leanname{eventually_prohibited_row_comparison},
\leanname{eventually_degree_cost_comparison},
\leanname{eventually_integer_degree_cost},
\leanname{eventually_integer_degree_sum},
\leanname{eventually_variable_tuple_degree_deletion},
\leanname{eventually_state_cost_comparison},
\leanname{eventually_density_cost_comparison}, via the state-cost comparison,
and \leanname{eventually_actual_padding_weight}, the last via
\leanname{eventually_prime_state_weight}, each retain the exponent of
their predecessor. The theorem
\leanname{eventually_padding_atom_comparison}, whose exponents $A_1,A_2$
come from \leanname{eventually_degree_cost_comparison} and
\leanname{eventually_density_cost_comparison}, and hence
\leanname{eventually_variable_padding_comparison}, takes
$A_1+A_2=2\cdot144789=289578$.
All are at most $10^6$. Enlarging $A$ strengthens the hypothesis
$U\ge\exp(L^A/2)$, proving the stated common choice.
\end{proof}

\section{Positive budgets and a rough-number multiplier}

Fix, for the rest of Part I,
\begin{equation}\label{eq:fixedparameters}
 A=1000000,\qquad W=10^{180},\qquad
 \delta=10^{-5},\quad\gamma=1-2\delta,\quad s_0=10^{-5}.
\end{equation}
These constants precede every sufficiently-large threshold.
For an integer $N$ put
\begin{equation}\label{eq:variableparameters}
 L=(\log N)^{1/A},\quad
 J=\left\lfloor\frac{\delta\log L}{6W}\right\rfloor,
 \quad\eta=\exp(-J),\quad K=\exp(4J),\quad
 Y_i=L^{1-\delta}\exp(6Wi).
\end{equation}
We henceforth take $N$ large enough that $J\ge1$.
For $1\le i\le J$ let
\[
 \mathcal P_i=\{p:\ p\text{ prime},\ p\equiv1\pmod5,
             \ p\nmid N,\ Y_{i-1}\le\log p<Y_i\},
 \qquad\mathcal P=\bigcup_i\mathcal P_i,
\]
and put
\[
 H_0=\exp(L^{1-\delta}),\quad
 \mathcal Q_0=\{p\le\exp L:\ p\text{ prime},\ p\not\equiv1\pmod5\},
 \qquad\mathcal Q=\{p\in\mathcal Q_0:p\nmid N\}.
\]
The fixed-modulus prime number theorem and partial summation give,
uniformly in $i$, the full band mass $3W/2+o(1)$.
Deleting prime factors of $N$ from all center bands costs at most
\begin{equation}\label{eq:centerholes}
 \sum_{p\mid N,\ p\ge H_0}\frac1p
 \le\frac{\log N}{H_0\log H_0}=o(1).
\end{equation}
Thus $W\le V_i\le2W$ for all $i$ once $N$ is large.
The estimate is uniform in $N$, because $\log N=L^A$ and
$H_0=\exp(L^{1-\delta})$.

Let $\mathcal D$ consist of products $d=\prod_{i=1}^J p_i$,
$p_i\in\mathcal P_i$, and let $V_*=\prod_iV_i$.
The bands are disjoint, so the tuple is determined by $d$, and
\begin{equation}\label{eq:tuplegeometry}
 \log d<\sum_{i=1}^J Y_i<2L,\qquad
 \sum_{d\in\mathcal D}\frac1d=V_*.
\end{equation}
Every pool prime is at most $\exp L$.
For squarefree products $q$ of primes in $\mathcal Q$ put
\[
 u(q)=4^{\omega(q)},\qquad
 w(x)=5^{\omega_{\mathcal Q}(x)},\qquad
 S=\prod_{p\in\mathcal Q}(1+4/p).
\]
For arbitrary integers $x$, $\omega_{\mathcal Q}(x)$ counts divisibility
by this finite prime set. In particular this definition also applies
at negative integers and at zero.

Retain pairs $(d,q)$ with $\omega(q)\le100\log L$ and
$\log(qd)\le100L+\eta$. Define disjoint bins
\begin{equation}\label{eq:bins}
 \mathcal S_j=\{(d,q)\text{ retained}:j\eta\le\log(qd)<(j+1)\eta\},
 \qquad D_j=\exp(j\eta),\quad T_j=ND_j.
\end{equation}
Only nonempty bins are used; their number is $O(L\exp J)$.
The union $\mathcal S$ of all bins is the word supply in
Definition~\ref{def:words}.
For each $d$ the retained divisor set is denoted by $\mathcal Q_d$.
The whole-product cutoff is included in $\mathcal Q_d$ before binning.

\begin{lemma}[Retained mass]\label{lem:mass}
For all sufficiently large $N$,
\begin{equation}\label{eq:mass}
 S_0:=\sum_j\sum_{(d,q)\in\mathcal S_j}\frac{u(q)}{qd}
 \ge\frac12SV_*.
\end{equation}
Moreover, $\#\mathcal S<\exp(101L)$ and
$m_Q:=\lfloor100\log L\rfloor$ is a valid padding-factor cap.
\end{lemma}
\begin{proof}
The probability law proportional to $u(q)/q$ independently selects each
padding prime with probability $4/(p+4)$.
Its mean $\log q$ is at most $4L+O(1)$, by
$\sum_{p\le\exp L}\log p/p=L+O(1)$.
For every $d$, \eqref{eq:tuplegeometry} and Markov give
$\Prob(\log(qd)>100L+\eta)<1/10$ for large $L$.
Also
\[
 \E\exp(\omega(q))
 \le\exp\bigl(4(\exp(1)-1)\sum_{p\le\exp L}1/p\bigr)
 \ll L^{4(\exp(1)-1)},
\]
so $\Prob(\omega(q)>100\log L)=o(1)$.
At least half of the padding mass is retained for each $d$; summing
$1/d$ proves \eqref{eq:mass}. These upper-tail estimates remain valid
after deleting any prime factors of $N$.
Finally, $dq$ determines its P-part $d$ and Q-part $q$, and every retained
product is at most $\exp(100L+\eta)<\exp(101L)$.
\end{proof}

For $d\in\mathcal D$ define the positive center envelope
\[
 a_d^+(x)=\prod_{p\mid d}(\ind_{p\mid x}+1/p),\qquad
 \rho_{j,d}(x)=\frac1{w(x)}
       \sum_{q:\,(d,q)\in\mathcal S_j}u(q)\ind_{q\mid x}.
\]
The densities sum to at most 1 over all bins, since
$\sum_{q\mid x,\ q\text{ squarefree on }\mathcal Q}u(q)=w(x)$.
Let
\[
 F_d(x)=w(x)\sum_j\rho_{j,d}(x)
      \ind_{\{\rho_{j,d}(x)>K/L\ \text{or}\
                         \omega_{\mathcal Q}(x)>400\log L\}}.
\]
The corresponding definitions for $\mathcal Q_0$ use the same
$\mathcal Q_d$ and bins, replacing only the vertex weight and degree
by their full-pool versions.

\begin{inputthm}[Full-pool padding cost]\label{input:paddingcost}
There is an absolute $C_p>0$ such that, for all sufficiently large $L$,
for every prime family with Q-pool $\mathcal Q_0$, every retained
squarefree divisor subset, every finite bin set, every
$0<\eta\le1$, every $K>0$, and every integer offset,
\begin{equation}\label{eq:fullpadding}
 \frac{\E\{a_d^+(x)F_{d,0}(x)\}}{S_{\rm full}}
 \le\frac{2^J}{d}\{C_p/K+L^{-100}\},
 \qquad S_{\rm full}=\prod_{p\in\mathcal Q_0}(1+4/p).
\end{equation}
The center primes in $d$ belong to the disjoint P-pool.
\end{inputthm}
This is the squarefree-tuple specialization of
\path{ModFiveThetaInput.eventually_positive_padding_cut_cost}, with
its fixed excluded set $E=\varnothing$; the proved input is
\path{modFiveThetaInput}.
The hypothesis on its divisor subset is just containment in the
squarefree products of the Q-pool. In particular it permits
$\mathcal Q_d$ rather than a complete divisor set.
In the source the excluded set is fixed before the threshold in $L$;
all the displayed prime families, subsets, bins and offsets follow
that threshold.

\begin{lemma}[Padding holes]\label{lem:paddingholes}
Uniformly in $N$ and in the bins and retained sets just defined,
\begin{equation}\label{eq:paddingholes}
 \frac1{SV_*}\sum_{d\in\mathcal D}\E\{a_d^+(x)F_d(x)\}
 \ll 2^J\{K^{-1}(\log L)^{10}+L^{-90}\}.
\end{equation}
\end{lemma}
\begin{proof}
By translation invariance of the independent uniform residue law, take
the single site to be zero. At a fixed site $x$, the equivalent
condition is $R_p+x\ne0\pmod p$ for every removed prime.
Put $E=\mathcal Q_0\setminus\mathcal Q$ and let
$A_E=\{R_p\ne0\text{ for every }p\in E\}$ in the full product model.
Then $\Prob(A_E)=\prod_{p\in E}(1-1/p)$.
Conditioned on $A_E$, the remaining coordinates have exactly the
reduced-pool law; at this one site the full and reduced weights,
degrees and retained-bin densities agree. The retained set
$\mathcal Q_d$ is identical on the two sides.
The center envelope uses disjoint coordinates. Nonnegativity gives
\[
 \frac{\E\{a_d^+F_d\}}S
 \le\frac{\E\{a_d^+F_{d,0}\}}{S_{\rm full}}
           \prod_{p\in E}\frac{p+4}{p-1}.
\]
The sum of reciprocals of the prime divisors of $N$ is at most
$\log\log\log N+O(1)$: use Mertens up to $\log N$ and the bound
$\omega(N)\le\log N/\log2$ above it.
Consequently the last product is
$O((\log\log N)^{10})=O_A((\log L)^{10})$.
Apply \eqref{eq:fullpadding} and sum $1/d=V_*$.
The extra logarithmic factor on $L^{-100}$ is absorbed by $L^{-90}$.
Only single-site padding costs have been conditioned in this argument.
\end{proof}

\begin{lemma}[Center and prohibited costs]\label{lem:rare}
Uniformly in all the parameters above and in any residue translation,
the positive center-degree cost, summed over all bins, is at most
\begin{equation}\label{eq:centertail}
 S2^JV_*\exp(-2WJ)
\end{equation}
when vertices with $\omega_{\mathcal P}(x)>6WJ$ are deleted.
The positive prohibited-site cost is at most
\begin{equation}\label{eq:prohibitedcost}
 SV_*\exp(-L^{0.9}).
\end{equation}
These costs have the same bounds for integer intervals of the bin
lengths, up to an error $\exp(-\tfrac12L^9)$.
\end{lemma}
\begin{proof}
For a fixed $d$, the mean of $a_d^+$ is $2^J/d$.
Tilting by this envelope makes a selected tuple prime active with
probability $(p+1)/(2p)$; all other center primes retain probability
$1/p$. The coordinates remain independent.
Their exponential moment is bounded by
$\exp((\exp(1)-1)(2W+1)J)$.
Since $W\ge10$,
\[
 \Prob_{a_d^+}(\omega_{\mathcal P}>6WJ)
 \le\exp(((\exp(1)-1)(2W+1)-6W)J)
 \le\exp(-2WJ).
\]
Multiply by $2^J/d$, sum over $d$, and use independence of the Q-pool
and $\sum_j\rho_{j,d}\le1$. This gives \eqref{eq:centertail}.

For prohibited sites the positive row envelope is
\[
 A_+(x)=w(x)\prod_{i=1}^J
             (\omega_{\mathcal P_i}(x)+V_i).
\]
Mertens and independence give
\[
 \E w^2=\prod_{p\in\mathcal Q}(1+24/p)\ll L^{24},\qquad
 \E A_+^2\le C L^{24}\prod_i(V_i+4V_i^2)
                  \le C L^{24}5^J V_*^2.
\]
The prohibited probability is bounded by \eqref{eq:badprob}, with
$h=N$, $s=\lfloor L^{1/10}\rfloor$, and the full retained supply
$\mathcal S$. Its hypotheses hold because
$J+m_Q\le101\log L$, both pool masses are at most $L^2$, and
$H_0\ge\exp(L^{199/200})$.
Cauchy--Schwarz bounds the positive rare-row mean by
$C L^{12}5^{J/2}V_*\exp(-\tfrac14L^{199/200})$.
Since $J=O(\log L)$ and $S\ge1$, this is at most the right side of
\eqref{eq:prohibitedcost}, even after multiplying by the number of bins.

For the integer statements use the positive-atom comparisons in
Input~\ref{input:comparison}, individually before summing.
There are at most $\exp(O(L))$ atoms and polynomially many bins.
The selected-prime count is at most $L^2$, each selected padding divisor
has at most $100\log L$ prime factors, each whole divisor has at most
$J+100\log L\le L^2$ prime factors, the supply has at most
$\exp(101L)$ pairs, $W\ge10$ and
$6WJ\le\delta\log L\le400\log L$. These are the hypotheses of
\leanname{eventually_prohibited_row_comparison},
\leanname{eventually_variable_tuple_degree_deletion} and
\leanname{eventually_variable_padding_comparison}.
Every interval length $\lfloor T_j\rfloor$ is at least $N$, and hence
at least $\exp(L^A/2)$.
The total comparison error is thus $\exp(-\tfrac12L^9)$ for large $L$.
Arbitrary offsets are included in the comparison. Product-law
translation invariance supplies the target-site versions.
\end{proof}

\begin{lemma}[Rough multiplier]\label{lem:rough}
Fix $0<\gamma_1<\kappa_1<1$.
For all sufficiently large $L$, let $y=\exp(L^{\gamma_1})$ and
$D\ge\frac12\exp(L^{\kappa_1})$.
If $Z\subset[D,2D)\cap\Z$ consists of $y$-rough integers and
$|c_n|\le1$, then
\begin{equation}\label{eq:rough}
 \norm{\sum_{n\in Z}\frac{c_n}{n}\e(n\theta)}_\infty\ll L^{-\gamma_1},
 \qquad
 \int_0^1\abs{\sum_{n\in Z}\frac{c_n}{n}\e(n\theta)}^4d\theta
             \ll D^{-1}L^{-4\gamma_1}.
\end{equation}
The constants may be taken absolute; the threshold may depend on
$\gamma_1,\kappa_1$.
\end{lemma}
\begin{proof}
We bound both one rough integer and the four forms $a,b,c,a+b-c$ in
the three-dimensional box $[D,2D)^3$.
For a prime $p$ their simultaneous nonvanishing density is exactly
\[
 \frac{(p-1)^3-(p-1)(p-2)}{p^3}
 =1-4/p+6/p^2-3/p^3.
\]
Indeed there are $(p-1)(p-2)$ nonzero pairs $a,b$ with $a+b\ne0$,
each excluding one nonzero value $c=a+b$.
The bad-event density $\nu_p$ is at most $4/p$.
Take the even Bonferroni order $r_0=2\lceil500\log L\rceil$.
The independent truncation error is at most
\[
 \frac{(\sum_{p\le y}\nu_p)^{r_0+1}}{(r_0+1)!}
 \le\frac{(5\log L)^{r_0+1}}{(r_0+1)!}\ll L^{-100}.
\]
Each truncated subset has CRT modulus $q\le y^{r_0}$.
The count in a fixed residue triple is
$(D/q)^3+O((D/q)^2+D/q+1)$.
Summing over its at most $q^3$ residue triples gives normalized error
$O(q/D)$, since $q\le D$ for large $L$.
There are at most $(r_0+1)(y+1)^{r_0}$ subsets, so the total error is
\[
 D^{-1}\exp(O(L^{\gamma_1}\log L))\ll L^{-100}.
\]
The one-dimensional CRT estimate has the same or smaller error.
Mertens' theorem bounds the independent good densities by
$C L^{-\gamma_1}$ and $C L^{-4\gamma_1}$, respectively; the ratio
of the four-form product to $\prod_{p\le y}(1-1/p)^4$ is bounded
because the remaining logarithmic terms are $O(p^{-2})$.
Consequently
\[
 |Z|\ll D L^{-\gamma_1},\qquad
 \#\{a,b,c,d\in Z:a+b=c+d\}\ll D^3L^{-4\gamma_1}.
\]
The first inequality proves the sup norm bound.
Expanding the fourth moment and integrating leaves precisely the
additive quadruples; each coefficient has absolute value at most
$D^{-4}$. The second inequality gives the remaining assertion.
\end{proof}

\section{The reflection graph and proof of the main theorem}

Extend $\lambda$ by zero to nonpositive integers only in the test functions
of this section. For a bin set
\begin{align}
 R_j&=\frac1{T_j}\sum_{1\le x\le T_j}
       \sum_{(d,q)\in\mathcal S_j}u(q)\ind_{qd\mid x}
                                  \lambda(x)\lambda(Nqd-x),\label{eq:raw}\\
 C_j&=\frac1{T_j}\sum_{1\le x\le T_j}
       \sum_{(d,q)\in\mathcal S_j}u(q)\ind_{q\mid x}
       \prod_{p\mid d}(\ind_{p\mid x}-1/p)
                                  \lambda(x)\lambda(Nqd-x).
                                                    \label{eq:centered}
\end{align}
Every real cutoff in these sums means its integer floor; the displayed
normalizing denominator remains the real number $T_j$.

\begin{lemma}[Dilation identity]\label{lem:dilation}
\begin{equation}\label{eq:dilation}
 \sum_jR_j=S_0\frac{C(N)}N+O(SV_*(\eta+N^{-1})).
\end{equation}
\end{lemma}
\begin{proof}
For each pair in \eqref{eq:raw}, put $x=qd\,m$.
Complete multiplicativity cancels $\lambda(qd)^2=1$, and the cutoff
$T_j/(qd)$ lies in $(N\exp(-\eta),N]$.
The lost tail of a 1-bounded sum is at most $N\eta+1$.
Also $T_j/(Nqd)\in(\exp(-\eta),1]$.
Thus that summand equals
$u(q)(qd)^{-1}C(N)/N+O(u(q)(qd)^{-1}(\eta+N^{-1}))$.
Summing proves the assertion, since $S_0\le SV_*$.
No total $N$ has been averaged.
\end{proof}

\begin{lemma}[CRT transfer]\label{lem:crt}
Let $P_*=\prod_{p\in\mathcal P\cup\mathcal Q}p$.
Choose a positive integer $\ell$ with $N\ell\equiv1\pmod{P_*}$.
For every $r,k,\xi\in\Z$ and every pool prime,
\begin{equation}\label{eq:crt}
 p\mid Nk+r+N\xi\quad\Longleftrightarrow\quad
 p\mid k+\ell r+\xi.
\end{equation}
All word divisibility, centered factors, Q-weights and Q-degree tests
therefore transfer from physical multiplier $N$ to quotient multiplier
1. The full prohibited predicate satisfies
\begin{equation}\label{eq:crtprohibited}
 \mathrm{Prohibited}_N(Nk+r)
 \quad\Longleftrightarrow\quad
 \mathrm{Prohibited}_1(k+\ell r).
\end{equation}
\end{lemma}
\begin{proof}
Every pool prime is coprime to $N$, and multiplication by $\ell$
proves \eqref{eq:crt}. The squarefree whole divisors $qd$ consequently
have the same divisibility at every departure point of a numerical word.
Its displacements scale by $N$. The suffix condition in a
forward-prohibited word becomes
$p\mid N\Delta_1(w_{\ge a})\iff p\mid\Delta_1(w_{\ge a})$,
since its $p$ belongs to the center pool.
The tuple intervals, unequal successive labels, length bounds and supply
membership do not change. Hence the forward-prohibited condition agrees
for every word, as does its minimality over oriented contiguous subwords.
Together with positivity this proves \eqref{eq:crtprohibited}.
\end{proof}
The corresponding finite identity is also supplied by
\path{ReflectedLiouville.prohibited_CRT} in the ancillary file
\texttt{CRTTransport.lean}; its hypotheses are precisely the
coprimality of every $qd$ with $N$ and the displayed inverses at pool
primes. The Liouville arguments are still the physical integers.

Fix a bin and write each integer uniquely as $x=Nk+r$, $1\le r\le N$.
Let $M_B=\lceil\exp(103L)\rceil$ and take two copies of
$[-M_B,M_B]\cap\Z$.
For every $d\in\mathcal D$ connect the left $k$ to the right
$t=k-qd$ whenever $(d,q)\in\mathcal S_j$, and include the transpose
edge. Put
\[
 c_d(x)=\prod_{p\mid d}(\ind_{p\mid x}-1/p),\quad
 \tau_{j,d}(x)=
 \ind_{\{\rho_{j,d}(x)\le K/L,\ \omega_{\mathcal Q}(x)\le400\log L,
                              \ x\text{ nonprohibited at multiplier }N\}}.
\]
The left-to-right entry of the resulting self-adjoint matrix is
\begin{equation}\label{eq:matrix}
 (B_{d,r})_{k,t}=
 \frac{L u(q)\tau_{j,d}(Nk+r)\tau_{j,d}(Nt+r)}
            {\sqrt{w(Nk+r)w(Nt+r)}}
       \ind_{q\mid Nk+r}c_d(Nk+r),\qquad t=k-qd.
\end{equation}
Every other left-to-right entry is zero. Let $E_r$ be the diagonal
projection onto center degree at most $6WJ$ in both layers.
Its values use the physical sites $Nk+r$, including negative sites.

\begin{proposition}[Compression]\label{prop:compression}
With the absolute choice $C_3=3^{170}$, for all sufficiently large
$N$, all but at most $N\exp(-2\lfloor L\rfloor)$ values of $r$ satisfy
\begin{equation}\label{eq:compression}
 \norm{E_r\sum_{d\in\mathcal D}B_{d,r}E_r}
 \le K(C_3\sqrt W)^J.
\end{equation}
The exceptional set may depend on the bin.
\end{proposition}
\begin{proof}
Let $g_r(k)=\sqrt{w(Nk+r)}$ and $a_d(k)=|c_d(Nk+r)|$.
For a fixed label, different $q$ give different neighbors in each layer.
The center factor and the test $q\mid Nk+r$ agree at the two endpoints
because $qd$ divides their physical difference $Nqd$.
From the density cap,
\[
 \sum_t|(B_{d,r})_{k,t}|\frac{g_r(t)}{g_r(k)}
 \le L a_d(k)\rho_{j,d}(Nk+r)
             \ind_{\{\rho_{j,d}(Nk+r)\le K/L\}}
 \le K a_d(k).
\]
The same calculation holds for a row in the right layer.
Weighted Cauchy--Schwarz and self-adjointness give
\begin{equation}\label{eq:squarefunction}
 \norm{B_{d,r}v}^2
 \le K^2\sum_k a_d(k)^2|v(k)|^2,
 \qquad\norm{B_{d,r}}\le K.
\end{equation}
To see the first inequality, bound each row square by
$K a_d(k)\sum_t|B_{k,t}|g_r(k)g_r(t)^{-1}|v(t)|^2$,
sum in $k$, and use $a_d(k)=a_d(t)$ on every nonzero entry.
The finite lemma \path{ReflectedLiouville.two_layer_localized_square} in
the ancillary file \texttt{TwoLayerRows.lean} has exactly this weighted-row
specialization and permits the duplicated site projection.
On $E_rv=v$, the factorization over the disjoint bands gives
\[
 \sum_d a_d(k)^2
 \le\prod_i(\omega_{\mathcal P_i}(Nk+r)+V_i)
 \le(8W)^J.
\]
The last inequality is AM--GM, because the sum of these $J$ factors
is at most $6WJ+2WJ$.
Thus \eqref{eq:spectralinput} applies to the self-adjoint family
$(B_{d,r})_{d\in\mathcal D}$ and the projection $E_r$ with the
nonnegative reals $a=K$ and $b=K(8W)^{J/2}$, by
\eqref{eq:squarefunction} and the last display.

Let $\mathcal H_r$ be the nonbacktracking operator on the direct sum of
one copy of the two-layer space for each $d$:
$(\mathcal H_r v)_d=\sum_{e\ne d}B_{e,r}v_e$.
All $B_{e,r}$ are real symmetric, so
$(\mathcal H_r^*v)_d=B_{d,r}\sum_{e\ne d}v_e$,
$\norm{\mathcal H_r^k}_{\HS}=\norm{(\mathcal H_r^*)^k}_{\HS}$ and
$\rad(\mathcal H_r^*)=\rad(\mathcal H_r)$.
Set $k_*=\lfloor L\rfloor$ and $B_0=K(2\exp(150)\sqrt W)^J$,
and expand $\norm{(\mathcal H_r^*)^{k_*}}_{\HS}^2$.
At each step the factor is $B_{d,r}$, where $d$ is the tag of the
arrival state. This is the convention of the source's
\leanname{shiftMatrix} and \leanname{integerShiftNext}: the new state
is the pair consisting of the step tag and the shifted site.
An entry from $(e_0,x_0)$ to $(d_{k_*},x_{k_*})$ sums paths with labels
$d_1,\ldots,d_{k_*}$ satisfying $d_1\ne e_0$ and
$d_i\ne d_{i+1}$. The squared Hilbert--Schmidt norm sums pairs of
such paths with the same initial and terminal states.
The initial tag $e_0$ is the fixed root tag, counted among the
$\le\exp(106L)$ states below; the common terminal tag is the common
last step label $d_{k_*}=d'_{k_*}$.
Reversing the second path gives a closed numerical word of length
$2k_*$, nonbacktracking in each half, with unconstrained seams.
The pair belongs to \leanname{actualClosedPairCatalog} for
$(e_0,x_0)$, since \leanname{closedPairNumericalTest} compares the
full end states: last label and integer site.
For each numerical word and initial layer the lifted path is unique.
Thus there are at most two lifts of each projected word.
The layer and finite-block gates are deterministic and independent of
the residue variables. The scalar masked weight of each lift equals the
projected weight. Taking the expectation of each signed word first, then
its absolute value, therefore puts these words within the majorant of
Input~\ref{input:trace}.

Its hypotheses hold with $h=1$, $C=101$,
$H=\lceil H_0\rceil$, $Y=B=\lfloor\exp L\rfloor$,
$Q_{\max}=\lfloor\exp(100L+1)\rfloor$ and
$D_{\max}=\lfloor\exp(2L)\rfloor$.
In particular $H_0\ge\exp(L^{199/200})$, the partial tuple products
are at most $\exp(2L)$, and $J+m_Q\le101\log L$.
The nonbacktracking state count is at most
\[
 2(2M_B+1)\#\mathcal D\le\exp(106L).
\]
The at most $2^{2k_*}$ direction choices, and the factor two for the
lift, are absorbed by the $\exp(108L)$ margin in
\eqref{eq:wordmajorant}, since
$\exp(106L)2^{2k_*}\cdot2\le\exp(108L)$ for large $L$.

For the integer-origin average use Lemma~\ref{lem:crt}.
Increase $\ell$ by a multiple of $P_*$ so that $\ell>M_B$.
At any initial quotient coordinate $k_0\in[-M_B,M_B]$, the origins
$\ell r+k_0$, $1\le r\le N$, have nonnegative initial origin
$a=\ell+k_0$ and step $\ell$.
Input~\ref{input:comparison} applies with length $N=\exp(L^A)$;
there is no bound on $\ell$.
All actual-word hypotheses hold: $2k_*+1\le4L$, $2k_*J\le L^2$,
$s\le L$, $J+m_Q\le L^2$, and $\#\mathcal S\le\exp(101L)$.
Each step has at most $2\#\mathcal S\le\exp(102L)$ numerical choices.
The total word and state comparison error is at most
\[
 \exp(106L+204L^2-L^9)=o(1).
\]
Combining it with the product-model majorant yields
\[
 \frac1N\sum_r\norm{\mathcal H_r^{k_*}}_{\HS}^2
 \le3B_0^{2k_*}.
\]
For a finite matrix, $\rad(\mathcal H_r)^{2k_*}$ is at most this
squared Hilbert--Schmidt norm. Markov therefore shows that
$\rad(\mathcal H_r)>\exp(3)B_0$ on at most
$3N\exp(-6k_*)\le N\exp(-2k_*)$ origins.
At every other origin apply \eqref{eq:spectralinput} with $a=K$,
$b=K(8W)^{J/2}$ and $\rad(\mathcal H_r)\le\exp(3)B_0$.
Since $a,b\le B_0$, it gives
$3\exp(3)B_0=3\exp(3)K(2\exp(150)\sqrt W)^J$, which for $J\ge1$
is at most $K(6\exp(153)\sqrt W)^J\le K(3^{170}\sqrt W)^J$.
This is \eqref{eq:compression}.
\end{proof}

Let $C_j^\circ$ denote \eqref{eq:centered} after keeping both endpoints
according to $\tau_{j,d}$ and $E_r$.

\begin{proposition}[The kept and deleted bilinear forms]\label{prop:bilinear}
For all sufficiently large $N$,
\begin{align}
 |C_j^\circ|&\ll\frac{SK(C_3\sqrt W)^J}{L}+\exp(-L),
                                                 \label{eq:kept}\\
 \frac{\sum_j|C_j-C_j^\circ|}{SV_*}
 &\ll2^J\{K^{-1}(\log L)^{10}+L^{-90}+\exp(-2WJ)\}
          +\exp(-L^{0.9}).\label{eq:deleted}
\end{align}
\end{proposition}
\begin{proof}
At origin $r$ use the unweighted tests
\[
 F_r(k)=\lambda(Nk+r)\ind_{\{0<Nk+r\le T_j\}},\quad
 G_r(t)=\lambda(-Nt-r)\ind_{\{0<-Nt-r\le N\exp(\eta)D_j\}}.
\]
Multiply each by its square-root weight, Q-degree cut and center
projection; put $f_r$ in the left layer and $g_r$ in the right layer.
All their supports lie in the block because
$D_j\le\exp(100L+\eta)$.
Each has $O(D_j)$ supported coordinates.
Input~\ref{input:comparison}, applied to the truncated weights on the
positive and negative physical intervals, gives
\[
 \sum_r(\norm{f_r}^2+\norm{g_r}^2)\ll ND_jS.
\]
Cancellation of the square-root weights gives the exact identity
\begin{equation}\label{eq:pairing}
 \sum_r\left\langle f_r,
             \bigl(E_r\sum_dB_{d,r}E_r\bigr)g_r\right\rangle
       =LT_j C_j^\circ.
\end{equation}
Here $t=k-qd$ corresponds to the physical graph site
$Nt+r=x-Nqd$. The test reads $Nqd-x=-Nt-r$.
All deletion predicates continue to be read at $Nt+r$.

For good origins use \eqref{eq:compression} and Cauchy--Schwarz in $r$.
At a bad origin bound the pairing row by row on the support of $F_r$.
The weighted absolute row bound is at most
$K w(Nk+r)\sum_d|c_d(Nk+r)|$.
On the kept test support,
$w(Nk+r)\le L^{400\log5}$ and
$\sum_d|c_d(Nk+r)|\le(8W)^J$.
The bad-origin cost is at most
$CND_j\exp(-2k_*)K L^{400\log5}(8W)^J$.
After division by $LT_j=LND_j$ it is $O(\exp(-L))$.
This counts only $O(D_j)$ source points per bad origin, and proves
\eqref{eq:kept}.

To bound the deletions replace $|c_d(x)|$ by $a_d^+(x)$ and use
$|\lambda|\le1$. At the source endpoint, the padding/Q-degree,
center-degree and prohibited costs are respectively bounded by
Lemma~\ref{lem:paddingholes} and Lemma~\ref{lem:rare}.
For the target endpoint, first translate each individual $(d,q)$ atom
by $-Nqd$. The test $q\mid x$ and envelope $a_d^+(x)$ are unchanged
by that translation, since $qd$ divides the shift.
The resulting atom is the same positive source cost on a translated
interval. The finite comparisons permit this arbitrary offset.
Sum the two endpoint bounds only after these individual translations.
The real normalization $T_j$ differs from the integer length
$\lfloor T_j\rfloor\ge N$ by less than $1$. With $\exp(O(L))$ atoms,
each bounded by a fixed power of $L$, this adds at most
$\exp(O(L))/N\ll\exp(-L^9)$ in addition to the comparison errors.
This proves \eqref{eq:deleted}.
\end{proof}

\begin{proposition}[Centering]\label{prop:centering}
\begin{equation}\label{eq:centering}
 \frac{\sum_j|R_j-C_j|}{SV_*}
 \ll\exp(J)2^J L^{-s_0}.
\end{equation}
\end{proposition}
\begin{proof}
Expand the centered product in \eqref{eq:centered}.
For every nonempty subset of mean-prime bands, let $v$ be the product
of its selected primes and $d'$ the product from the other bands.
After putting $x=qd'm$, complete multiplicativity cancels
$\lambda(qd')^2$. For each fixed $q,d'$ and bin, the remaining sum,
apart from its external factor $u(q)/(qd')$, has the form
\begin{equation}\label{eq:convolution}
 \frac1Y\sum_{m\le Y}\sum_{v\in Z}\frac{c_v}{v}
                                \lambda(m)\lambda(Nv-m),\qquad |c_v|\le1.
\end{equation}
The bin restrictions are included in $Z$ and $c_v$.
At most two dyadic pieces suffice to arrange
$Z\subset[D,2D)$, $Y\asymp ND$, and
$H_0/2\le D\le\exp(2L)$.
Every $v$ is $\exp(L^\gamma)$-rough.
The factor outside \eqref{eq:convolution} is exact, because
$Y=T_j/(qd')$.

Write $m=Nk+r$ and define
\begin{align*}
 F_r(\xi)&=\sum_{k\ge0,\ Nk+r\le Y}\lambda(Nk+r)\e(k\xi),\\
 G_r(\xi)&=\sum_{1\le t\le\lceil2D\rceil}\lambda(Nt-r)\e(t\xi),\\
 B(\xi)&=\sum_{v\in Z}(c_v/v)\e(-v\xi).
\end{align*}
The integral of $F_rG_rB$ forces $k+t=v$, so
\eqref{eq:convolution} is exactly
$Y^{-1}\sum_r\int_0^1F_rG_rB$.
One has $|G_r|\ll D$ and
$\norm{F_r}_2\norm{G_r}_2\ll D$ by Parseval.
Lemma~\ref{lem:fourier} with $\nu_0=\gamma/A$ gives
\[
 \frac1N\sum_r|F_r(\xi)|\ll D L^{-1/4000}.
\]
The prefix lengths for different $r$ differ by at most one term; a
single common prefix and that error justify this application.
Its lower-size condition follows from
$\log D\ge\tfrac12L^{1-\delta}\ge L^\gamma$ for large $L$, and
$(1-\delta)/3000>1/4000$.
Also $D\le\exp(2L)\le N$.
Lemma~\ref{lem:rough}, with $\gamma_1=\gamma$,
$\kappa_1=1-\delta$, gives
\[
 \norm B_\infty\ll L^{-\gamma},\qquad
 \int_0^1|B|^4\ll D^{-1}L^{-4\gamma}.
\]
On $|B|\le L^{-1-s_0}$, Parseval bounds the convolution by
$O(L^{-1-s_0})$.
The complementary set has measure at most
$C D^{-1}L^{4(1+s_0)-4\gamma}$.
There use the $r$-averaged $F_r$ bound, $|G_r|\ll D$ and the sup norm
of $B$. Since $Y\asymp ND$, its contribution is
\[
 O\bigl(L^{4(1+s_0)-5\gamma-1/4000}\bigr)
       =O(L^{-1-s_0}),
\]
because
\begin{equation}\label{eq:margin}
 5\gamma+1/4000-(5+5s_0)=1/10000>0.
\end{equation}
For each mean-prime subset, the sum of the external positive factors
$u(q)/(qd')$ is at most $S\prod_{i\notin I}V_i\le SV_*$.
There are fewer than $2^J$ such subsets and $O(L\exp J)$ bins.
Multiply these counts by $L^{-1-s_0}$ to obtain \eqref{eq:centering}.
\end{proof}

\begin{proof}[Proof of Theorem~\ref{thm:main}]
Sum \eqref{eq:kept} over the $O(L\exp J)$ bins and divide by $SV_*$.
Since $V_*\ge W^J$, the spectral contribution is at most
\[
 C\exp(5J)(C_3/\sqrt W)^J\le C\exp(-J),
\]
since $\exp(6)C_3<3^{176}<10^{88}<\sqrt W$.
All terms of \eqref{eq:deleted} are $O(\exp(-J))$:
\[
 2^JK^{-1}(\log L)^{10}=\exp(-(4-\log2)J)(\log L)^{10},
\]
and $J=\delta\log L/(6W)+O(1)$ with its positive coefficient fixed.
The terms $2^JL^{-90}$, $2^J\exp(-2WJ)$ and $\exp(-L^{0.9})$ are
smaller. Moreover
\[
 \exp(J)2^JL^{-s_0}=O(\exp(-J)),
 \qquad \frac{(2+\log2)\delta}{6W}<s_0,
\]
so \eqref{eq:centering} has the same bound.
Use \eqref{eq:mass} and \eqref{eq:dilation} to conclude
\[
 \frac{|C(N)|}{N}\ll\exp(-J)
 \le\exp(1)(\log N)^{-c_*},\qquad
 c_*:=\frac{\delta}{6WA}.
\]
The constants $A,W$ are exactly \eqref{eq:fixedparameters}.
With the numerical choices in \eqref{eq:fixedparameters},
\[
 c_*=(6\cdot10^{191})^{-1}>10^{-192}>10^{-200}.
\]
This proves \eqref{eq:correlation}.
Lemma~\ref{lem:realmean}, with the modulus-one character, bounds the
linear sum in \eqref{eq:indicator} by
$O(N(\log N)^{-1/40})$.
Insert that estimate and the correlation bound in
Lemma~\ref{lem:indicators}; the difference between $(N-1)/4$ and $N/4$
is absorbed by the error. This proves \eqref{eq:main}.
Choose $N_0$ so large that, for $N\ge N_0$, one has $J\ge1$;
$L=(\log N)^{1/A}$ exceeds the thresholds of
Inputs~\ref{input:trace}--\ref{input:bad-spectral} and
\ref{input:paddingcost}, with $h=1$ or $h=N$ as indicated,
$C=101$, $W=10^{180}$ and $E=\varnothing$;
of Lemma~\ref{lem:rough}, with $\gamma_1=\gamma$ and
$\kappa_1=1-\delta$; and of the finitely many inequalities for large $L$.
Also require $N$ to exceed the thresholds of
Lemmas~\ref{lem:realmean} and \ref{lem:fourier}, the latter with
$\nu_0=\gamma/A$, and of Lemma~\ref{lem:variance} with $\nu=\nu_0/2$.
The input thresholds precede the quantification of $h$ and the prime
families. All thresholds depend only on $A,W,\delta,\gamma,s_0,C_3$;
\eqref{eq:centerholes} and \eqref{eq:paddingholes} are uniform in the
prime factorisation of $N$. Thus one absolute $N_0$ works for both
parities and all four signs. The condition $J\ge1$ alone requires
$\log\log N\ge6\cdot10^{191}$.
\end{proof}

\part{Consequences of the zero-free half-plane}
\section{Theorem substitutions}
Besides the Lean declaration transcribed in Input~\ref{input:zf},
this part uses Pintz's explicit formula \cite{PintzI}, the zero-density
estimate of Chen--Gupta--Li \cite{CGL}, a preprint at the time of writing,
Montgomery--Vaughan's Theorem~13.12 \cite{MV}, and Martin's
Proposition~3 \cite{Martin1997}. Their applications here are written
mathematics, independent of the graph argument in Part I.
The exponents are tracked explicitly; implied constants retain the
dependencies indicated by their subscripts.
Theorems~\ref{thm:goldbach} and \ref{thm:leastprime} need some fixed
zero-free half-plane $\Rea s>1-\delta_0$, with $\delta_0>0$;
the exponents $8$ in Theorems~\ref{thm:nonresidue} and
\ref{thm:primitive}, and consequently $24$ for primitive roots,
use the specific value $7/8$.

\subsection{The binary Goldbach exceptional set}
Let $E(X)$ count the even integers $4\le m\le X$ which are not a sum
of two primes.

\begin{theorem}\label{thm:goldbach}
For every $\varepsilon>0$,
\begin{equation}\label{eq:goldbach}
 E(X)\ll_\varepsilon X^{3/5+\varepsilon}.
\end{equation}
\end{theorem}

Here is the exact major-arc input used. In Pintz,
\emph{A new explicit formula in the additive theory of primes with
applications I}, Theorem 1 and equation (2.10) of version 1
\cite{PintzI}, use its notation
\[
 X_1=X^{1-\varepsilon_0},\quad
 S_X(\alpha)=\sum_{X_1<p\le X}\log p\,\e(p\alpha),\quad
 Q=X/P,
\]
with the disjoint major arcs $|\alpha-a/q|\le1/(qQ)$, $q\le P$,
$(a,q)=1$. Write
\[
 R_1(m)=\int_{\mathfrak M}S_X(\alpha)^2\e(-m\alpha)\,d\alpha.
\]
Let $\mathfrak m=[1/Q,1+1/Q)\setminus\mathfrak M$ and put
$R_2(m)=\int_{\mathfrak m}S_X(\alpha)^2\e(-m\alpha)\,d\alpha$.
Orthogonality gives
$R_1(m)+R_2(m)=\sum_{p+p'=m,\ p,p'\in(X_1,X]}\log p\log p'$.
We use \eqref{eq:pintzminor} for even $m\in[X/2,X]$.
For a small fixed $\eta>0$, $P_0\le X^{4/9-\eta}$, a sufficiently
large fixed $H$ and sufficiently large $X$, it gives a single
$P\in[P_0X^{-\eta},P_0]$ such that, for all $m\le X$,
\begin{equation}\label{eq:pintz}
 \begin{split}
 R_1(m)={}&\sum_{(\rho_i,\chi_i),(\rho_j,\chi_j)\in\mathcal E}
 A(\rho_i)A(\rho_j)\mathfrak S(\chi_i,\chi_j,m)
 \frac{\Gamma(\rho_i)\Gamma(\rho_j)}{\Gamma(\rho_i+\rho_j)}
 m^{\rho_i+\rho_j-1}\\
 &+O_\eta(\mathfrak S(m)X\exp(-c_\eta H))
       +O_\eta(X^{1-\varepsilon_0}).
 \end{split}
\end{equation}
The set $\mathcal E$ in its equation (2.8) consists of the pole
$(1,\chi_0\pmod1)$ with $A(1)=1$, and nontrivial zeros of primitive
characters of conductor at most $P$, with
$\beta\ge1-H/\log X$, $|\gamma|\le\sqrt X$, and $A(\rho)=-1$.
All small constants, including $\varepsilon_0>0$, are fixed according
to the major-arc range; the allowed major-arc parameters also satisfy
$(\log X)^C\le P\le\sqrt X$.
The minor-arc input is \cite[(5.2)]{PintzI}, obtained from Parseval's
identity and Vaughan's estimate:
\begin{equation}\label{eq:pintzminor}
 |R_2(m)|\le X/\sqrt{\log X}
\end{equation}
for all $m\le X$ except
$\ll(\log X)^{10}\max(X/P,X^{3/5})$ values.
For $P\le X^{2/5}$ this is the estimate quoted as (2.6) in
\cite{PintzII}.

\begin{proof}
Fix $0<\eta<\min(\varepsilon/2,2/45)$ and take $P_0=X^{2/5}$.
Choose $H\ge H_0(\eta)$ first so that the first implied constant in \eqref{eq:pintz}
times $\exp(-c_\eta H)$ is less than $1/8$.
Then take $X$ sufficiently large for the theorem and $\log X>8H$.
Input~\ref{input:zf} leaves only the pole in $\mathcal E$, and hence
\[
 R_1(m)=\mathfrak S(m)m
   +O_\eta(\mathfrak S(m)X\exp(-c_\eta H))
   +O_\eta(X^{1-\varepsilon_0}).
\]
For even $m$,
$\mathfrak S(m)\ge2\prod_{p\ge3}(1-(p-1)^{-2})>0$.
Thus $\Rea R_1(m)\ge\mathfrak S(m)X/4\gg X$ for $X/2\le m\le X$.
Outside the exceptional set in \eqref{eq:pintzminor},
$R_1(m)+R_2(m)>0$, which counts two actual primes with positive
logarithmic weights. Since $P\ge X^{2/5-\eta}$, the exceptions in this
dyadic interval number at most
$C_\eta X^{3/5+\eta}(\log X)^{10}$.
Absorb the logarithms into $X^{\varepsilon-\eta}$ and sum the dyadic
intervals to prove \eqref{eq:goldbach}.
\end{proof}

The exponent $3/5$ is Pintz's: \cite[Theorem~2, (2.13)--(2.14)]{PintzI}
obtains it under the stated Siegel-zero hypothesis, and
\cite[Theorem~3, (2.16)]{PintzI} under a zero-free box adjacent to
$\sigma=1$. Its value comes from the minor-arc estimate
\cite[\S5]{PintzI}. Input~\ref{input:zf} makes this conditional exponent
unconditional by leaving only the pole in \eqref{eq:pintz}.
The argument uses only a fixed zero-free half-plane.
For the preceding unconditional bounds, Pintz
\cite[Theorem~1]{PintzII} obtained $0.72$, and Zhao's preprint
\cite[Theorem~1.1, version~2]{Zhao2026} states $0.7$.
The generalized-zero cutoff in \eqref{eq:pintz} is scaled by $\log X$.

\subsection{The least prime in a reduced progression}
For $(a,q)=1$ let $p(a,q)$ be the least prime congruent to $a\pmod q$.

\begin{theorem}\label{thm:leastprime}
For every $\varepsilon>0$, uniformly for positive integers $q$ and all
reduced residue classes $a\pmod q$,
\begin{equation}\label{eq:leastprime}
 p(a,q)\ll_\varepsilon q^{7/3+\varepsilon}.
\end{equation}
\end{theorem}

The proof uses Proposition~\ref{prop:ap-smooth}, which derives the needed
case directly from Chen--Gupta--Li's Theorem~1.2.
For comparison, Corollary~1.4, general-modulus clause in version~2 of
\emph{Large Value Estimates for Dirichlet Polynomials with Characters
and Zero Density of Dirichlet $L$-Functions} \cite{CGL}, assumes
$h\le x$, $q_1\mid q$, $q_1\ge\sqrt q$, $(a,q)=1$, and a common zero-free
region $\sigma\ge1-\eta(q,T)$, $|\Im s|\le T$ for
$\prod_{\chi\bmod q}L(s,\chi)$, with
\[
 \eta(q,x^{\epsilon/3}T_0)\log x\longrightarrow\infty
 \quad\text{uniformly for }q<h/x^{17/30-\epsilon},\qquad
 T_0=(x/h)^{1+o(1)}.
\]
With the usual meromorphic interpretation at the principal pole, its
conclusion is
\[
 \psi(x+h;q,a)-\psi(x;q,a)=(1+o_{\eta,\epsilon}(1))h/\ph(q)
\]
provided
\[
 h/\ph(q)\ge x^{1/2+\epsilon}q^{1/6}+x^{17/30+\epsilon}.
\]
The paper states this corollary with a reference to a smoothed explicit
formula. Proposition~\ref{prop:ap-smooth} supplies that deduction from
Theorem~1.2 for the case used here.
Input~\ref{input:zf} supplies its zero-free hypothesis at every modulus
and height with $\eta(q,T)=1/9$; thus
$\eta(q,x^{\epsilon/3}T_0)\log x\to\infty$ uniformly in $q$.
\begin{proof}
Given $\varepsilon>0$, take $\delta=\min(1/200,\varepsilon/100)$
and fix a nonzero nonnegative $\Phi\in C_c^\infty((1,2))$.
For $q\le x^{3/7-\delta}$, Proposition~\ref{prop:ap-smooth} gives
a prime $x<p<2x$ in every reduced class for $x\ge x_0(\delta,\Phi)$.
Set $x=q^{1/(3/7-\delta)}$. Since
\[
 \frac1{3/7-\delta}-\frac73
 =\frac{49\delta}{9-21\delta}\le6\delta<\varepsilon,
\]
we obtain \eqref{eq:leastprime}; the finitely many smaller moduli,
including $q=1$, are absorbed in its constant.
\end{proof}

The previous published exponent is $5$, due to Xylouris
\cite[Theorem~2.1]{Xylouris2011}, with an effective constant;
\cite[Theorem~1]{Xylouris2018} is a published condensed version
and states $p(a,q)<Cq^5$. Zhao \cite[Theorem~1.4]{Zhao2026} also
states $5$. Naslund's computer-assisted preprint \cite{Naslund2026}
states $p(a,q)<q^{3.99}$ for all sufficiently large $q$, uniformly in $a$.

\subsection{The least quadratic non-residue}
For an odd prime $p$ let $n(p)$ be its least positive quadratic
non-residue.

\begin{theorem}\label{thm:nonresidue}
Uniformly for all odd primes,
\begin{equation}\label{eq:nonresidue}
 n(p)\ll(\log p)^8.
\end{equation}
\end{theorem}

The input is Montgomery--Vaughan, \emph{Multiplicative Number Theory I},
Theorem 13.12 \cite{MV}, pages 428--430.
For a nonprincipal character $\chi\pmod q$ it assumes
$1/\log q\le\delta\le1/2$ and
\[
 L(s,\chi)\ne0\quad\text{for }1-\delta<\sigma<1,
                      \quad0<|t|\le\delta^2\log q.
\]
Its conclusion is
$n(\chi)<(A_{\rm MV}\delta\log q)^{1/\delta}$ for an absolute constant
$A_{\rm MV}$, where $n(\chi)$ is the least $n$ with
$\chi(n)\notin\{0,1\}$.
\begin{proof}
Take $\chi$ to be the Legendre character modulo $p$ and
$\delta=1/8$. For $p\ge\exp(8)$ the parameter range holds, and
Input~\ref{input:zf} supplies the stated zero-free rectangle.
Thus $n(p)<(A_{\rm MV}/8)^8(\log p)^8$; the remaining odd primes
are covered by $n(p)<p$ and an absolute enlargement of the constant.
\end{proof}
The classical unconditional comparison is Burgess
\cite[Theorem 2]{Burgess1957}:
$n(p)\ll_\varepsilon p^{1/(4\sqrt{\exp(1)})+\varepsilon}$.
OpenAI's accompanying manuscript already gives
$n(p)\le C(\log p)^A$, with $A=32$ as an example in
\cite[Corollary~1.2]{OpenAIQ78}, via Bhargava--Ivanyos--Mittal--Saxena.
The companion \cite{OpenAIQ1112} uses Montgomery--Vaughan's
Theorem~13.12 for a fixed power of $\log p$.
Equation~\eqref{eq:nonresidue} gives its explicit exponent
$1/\delta=8$ at $\delta=1/8$.


\subsection{Primitive roots and generators}
For a modulus $q$ whose unit group is cyclic, let $g(q)$ and $g^*(q)$
be the least positive integer and least prime primitive roots, respectively.

\begin{theorem}\label{thm:primitive}
For every odd prime $p$, with $r=\omega(p-1)$,
\begin{equation}\label{eq:primitive}
 g(p)\le g^*(p)\ll
 r^{16}(1+\log r)^{16}(\log p)^8\ll(\log p)^{24}.
\end{equation}
For every $q\ge3$ with cyclic unit group the same assertions hold with
$r=\omega(\ph(q))$ and $p$ replaced by $q$, with absolute constants.
These moduli are $4$, odd prime powers, and twice odd prime powers.
For $q=2$ one has $g(2)=1$ and $g^*(2)=3$.
\end{theorem}

The input is Martin, \emph{The least prime primitive root and the
shifted sieve}, Proposition~3 \cite{Martin1997}.
Write $E(q)$ for the exponent of the unit group, $s(m)$ for the
largest squarefree divisor of $m$, and
$\Phi^*(q)=\{\chi^{E(q)/s(\ph(q))}:\chi\pmod q\}$.
Let $c_0$ be the proportion of unit classes having order $E(q)$ and
$\log_1 u=\max(\log u,1)$.
The printed conditions of Proposition~3 are $q\ge2$ and
$1/2\le\sigma<1$. We use it for $q\ge3$, so
$\omega(\ph(q))\ge1$, and put
\[
 f(q,\sigma)=
 \{\omega(\ph(q))^2\log_1\omega(\ph(q))\,c_0^{-1}\log q\}^{1/(1-\sigma)}.
\]
Martin's hypothesis is that every nonprincipal $\chi\in\Phi^*(q)$
has no zero with $\Rea s>\sigma$ and $|\Im s|<f(q,\sigma)$.
His conclusion is $g^*(q)\ll_\sigma f(q,\sigma)$, where for general
$q$ his $g^*(q)$ denotes the least prime of maximal order.
His Lemma~6 gives the smoothed estimate
\[
 \sum_{n<x}\Lambda(n)\chi(n)(1-n/x)\ll x^\sigma\log q
\]
for those nonprincipal characters when $1\le x\ll T\ll q$ and
their $L$-functions have no zero with $\Rea s>\sigma$, $|\Im s|<T$.
At the polylogarithmic search length below one may take $T=x\ll q$.

\begin{proof}
Input~\ref{input:zf} supplies Martin's hypothesis with $\sigma=7/8$
at every height. For a cyclic group,
$c_0=\ph(\ph(q))/\ph(q)$. Put $n=\ph(q)$ and let $p_i$ be the
$i$th prime. Mertens' theorem and $p_r\ll r\log(2r)$ give
\[
 c_0^{-1}=\frac n{\ph(n)}
 \le\prod_{1\le i\le\omega(n)}(1-1/p_i)^{-1}
 \ll\log_1\omega(n),
\]
as in Martin's remark after Proposition~3. Thus
$f(q,7/8)\ll r^{16}(1+\log r)^{16}(\log q)^8$.
If $r\ge2$, the product of $r$ distinct prime divisors of $\ph(q)$
is at least $r!$, so $r(1+\log r)\ll\log q$ by the factorial bound;
the case $r=1$ follows since $q\ge3$.
Substitution yields $(\log q)^{24}$, and $g(q)\le g^*(q)$ is immediate.
\end{proof}

Corollary~\ref{cor:generators} specializes Bach's character and
proper-subgroup bound \cite[p.~293, unnumbered theorem and the subsequent
subgroup reduction]{Bach1982} to the uniform zero-free boundary $7/8$.

\begin{corollary}\label{cor:generators}
There is an absolute $C>0$ such that for every $q\ge2$ and every
nonprincipal character $\chi\pmod q$ there is a prime
$\ell\nmid q$, $\ell<C(\log q)^8$, with $\chi(\ell)\ne1$.
The unit group modulo $q$ is generated by the primes in this range.
For every $k\ge2$ whose $k$th-power subgroup is proper, the least
positive unit outside it is prime and less than $C(\log q)^8$,
with the same constant independent of $k$.
\end{corollary}
\begin{proof}
Apply Montgomery--Vaughan's Theorem~13.12, as stated above, with
$\delta=1/8$ to any nonprincipal $\chi$. It gives
$n(\chi)<C(\log q)^8$ for $q\ge\exp(8)$; finitely many smaller
moduli are included by enlarging $C$.
The least such unit $n(\chi)$ is prime: if it were composite,
all its prime factors would be smaller units on which $\chi$ equals $1$.
If the indicated primes generated a proper subgroup, the quotient
would have a nontrivial character, contradicting this conclusion.
A proper $k$th-power subgroup also has a nontrivial annihilating
character. Hence a prime in the stated range lies outside it;
the least positive unit outside a subgroup is prime by the same
factorisation argument.
\end{proof}

Sandlund \cite{Sandlund2026}, in his research memo dated 7 October 2026,
also derives proper-subgroup bounds and, via Shoup's sieve, least prime
primitive-root bounds from a fixed Dirichlet zero-free half-plane, and
constructs primitive roots when the factorisation of $p-1$ is supplied.

The classical unconditional bound
$g(p)\ll_\varepsilon p^{1/4+\varepsilon}$ is due to Wang in 1959,
using Burgess's character-sum estimate, as credited in
\cite[\S1]{Bach1997} and \cite[\S1]{FGG2026}; the same exponent
follows from Burgess \cite[Theorem~1]{Burgess1962}.
Indeed that theorem gives a power saving for nonprincipal character
sums of length $H=\lceil p^{1/4+\varepsilon}\rceil$ on choosing its integer parameter
larger than $1/(4\varepsilon)$; the primitive-root character expansion
then uses $2^{\omega(p-1)}=p^{o(1)}$ to give the stated bound.
Carella \cite[Theorem~1.1]{Carella2025} claims the unconditional bound
$g(p)\ll_\varepsilon(\log p)^{1+\varepsilon}$ for fixed small
$\varepsilon>0$ and sufficiently large primes.
For the least prime primitive root, Ha's publisher abstract
\cite{Ha2013} reports the unconditional exponent $3.1$.
Martin \cite[Theorem~1]{Martin1997} obtains a fixed power of $\log p$
for all but $O_{\varepsilon,\eta}(Y^{\varepsilon})$ odd prime powers
$p^t\le Y$, with $0<\varepsilon\le20/21$ and $\eta>0$.
Under GRH, Wang's bound, as recorded in \cite[\S1]{Shoup1992}, is
$g(p)\ll r^6(\log p)^2$. Shoup \cite[Theorem~3]{Shoup1992} obtains
$g(p)\ll r^4(1+\log r)^4(\log p)^2\ll(\log p)^6$;
Martin \cite[Corollary~3.1]{Martin1997} states the prime version and
extends it to composite moduli as an estimate for primes of maximal order.
Bach \cite[Theorem~3.2]{Bach1997}, under GRH and for odd primes $p$,
constructs a set of $O((\log p)^4/(\log\log p)^3)$ residue classes
containing a primitive root, using
$O((\log p)^7/(\log\log p)^3)$ bit operations with ordinary arithmetic.
The bound is for the number of candidate classes.

\appendix
\section{A uniform density deduction for progressions}
\begin{proposition}\label{prop:ap-smooth}
Fix $0<\delta<1/100$ and a nonzero nonnegative
$\Phi\in C_c^\infty((1,2))$. Uniformly for
$q\le x^{3/7-\delta}$ and $(a,q)=1$, as $x\to\infty$,
\[
 \sum_{n\equiv a\ (q)}\Lambda(n)\Phi(n/x)
   =(c_\Phi+o_{\delta,\Phi}(1))\frac{x}{\ph(q)},\qquad
 c_\Phi=\int_1^2\Phi(u)\,du>0.
\]
Every such class contains a prime at most $2x$ once $x$ is sufficiently
large in terms of $\delta,\Phi$.
\end{proposition}
The input is Chen--Gupta--Li,
\emph{Large Value Estimates for Dirichlet Polynomials with Characters
and Zero Density of Dirichlet $L$-Functions}, Theorem 1.2, last
universal estimate in version 2 \cite{CGL}.
For $q\ge1$ and $1/2<\sigma<1$, with $N(\sigma,T,\chi)$ counting
zeros $\beta+i\gamma$ with $\sigma\le\beta\le1$, $|\gamma|\le T$,
the final ``in all cases'' display of Theorem~1.2 states, with
the paper's convention $qT\to\infty$ for $o(1)$
(we use $T=x^\tau\to\infty$),
\begin{equation}\label{eq:cgl}
 \sum_{\chi\bmod q}N(\sigma,T,\chi)
 \le(qT)^{o(1)}
  \{(q^{7/3}T^2)^{1-\sigma}+(qT)^{(30/13)(1-\sigma)}\}.
\end{equation}
This clause has no smoothness condition on $q$.
Its $o(1)$ means that any fixed positive power of $qT$ is an admissible
loss with a corresponding fixed implied constant.
\begin{proof}
Fix $0<\delta<1/100$ and suppose $q\le x^{3/7-\delta}$.
Take $\tau=\delta/8$, $T=x^\tau$, and density loss
$\alpha=\delta/100$ in \eqref{eq:cgl}.
Then
\[
 q^{7/3}T^2\le x^{1-25\delta/12}\le x^{1-\delta},\qquad
 (qT)^{30/13}\le
 x^{90/91-105\delta/52}\le x^{1-\delta}.
\]
Fix a nonzero nonnegative $\Phi\in C_c^\infty((1,2))$, with Mellin
transform $\widetilde\Phi$ and $c_\Phi=\widetilde\Phi(1)>0$.
The smoothed explicit formula, from
\cite[Theorem 12.10 and the zeta formula in Chapter 12]{MV}, and
character orthogonality give
\[
 \sum_{n\equiv a\ (q)}\Lambda(n)\Phi(n/x)
 =\frac{c_\Phi x}{\ph(q)}
 +O_\Phi\left(\frac{x}{\ph(q)}
     \sum_{\chi\bmod q}\sum_\rho x^{\beta-1}
                   |\widetilde\Phi(\rho)|+\log^2(qx)\right).
\]
For primitive $\chi$, \cite[Theorem~12.10]{MV} and its zeta analogue,
integrated by parts against $\Phi'(u/x)/x$, give
\[
 \sum_n\Lambda(n)\chi(n)\Phi(n/x)
 =\ind_{\chi=\chi_0}\widetilde\Phi(1)x
   -\sum_\rho x^\rho\widetilde\Phi(\rho)+O_\Phi(1).
\]
The constant $C(\chi)$ in that theorem's (12.6) is independent of
$u$ and integrates to zero against $\Phi'$.
For fixed $x,q$, its remainder tends to zero after integration as the
truncation height tends to infinity: the printed remainder bound
gives pointwise convergence away from integers and an integrable
bound on this compact interval. The zero series converges absolutely
by $\widetilde\Phi(\beta+i\gamma)\ll_{\Phi,B}(1+|\gamma|)^{-B}$.
The induced Euler factors cost $O_\Phi(\log^2(qx))$ after orthogonality;
the principal pole supplies the main term.

All nontrivial zeros have $\beta\le7/8$.
For $|\gamma|\le T$ and $\beta\le1/2+\alpha$, the unit-height count
\cite[Theorem 10.17]{MV} bounds the normalized zero cost by
$O_\Phi(qT\log(qT)x^{-1/2+\alpha})=o(1)$.
Partition $[1/2+\alpha,7/8]$ into intervals of width at most $\alpha$.
For each lower endpoint $\sigma$, \eqref{eq:cgl} bounds its cost by
\[
 C_{\delta,\Phi}x^{\sigma+\alpha-1}
       x^\alpha x^{(1-\delta)(1-\sigma)}
 \le C_{\delta,\Phi}x^{-\delta/8+2\alpha}=o(1).
\]
This uses only finitely many fixed values of $\sigma$.
For $|\gamma|>T$, choose $B=\lceil10/\tau\rceil$ in the rapid decay
$\widetilde\Phi(\beta+i\gamma)\ll_{\Phi,B}(1+|\gamma|)^{-B}$.
The aggregate tail is
$O_{\delta,\Phi}(q x^{-1/8}T^{1-B}\log(qT))=o(1)$.
Thus the weighted progression sum is
$(c_\Phi+o(1))x/\ph(q)$, uniformly in the stated range and in $a$.
Prime powers of exponent at least two cost
$O(\sqrt x\log^2x)=o(x/\ph(q))$, since
$x/\ph(q)\ge x/q\ge x^{4/7+\delta}$.
There is therefore a prime at most $2x$ in every such class.
This proves the uniform asymptotic and the prime assertion.
\end{proof}
The same progression range follows from any fixed zero-free half-plane
$\Rea s>1-\delta_0$, with $0<\delta_0\le1/2$.
Take $\alpha=\min(\delta/100,\delta\delta_0/10)$, partition up to
$1-\delta_0$, and replace the grid cost by
$x^{-\delta\delta_0+2\alpha}=o(1)$ and the tail factor by
$x^{-\delta_0}$. The other bounds above are unchanged.


\bibliographystyle{plainnat}
\bibliography{references}
\section*{Use of AI}
The proofs and the text of this paper were produced with AI systems (GPT-6 Astra, GPT-6.1 Sol and Claude Opus~5.5) under the author's direction. They were checked by independent reviews with AI models; the inputs from the OpenAI library are Lean theorems, and the finite lemmas in the ancillary files are checked in Lean. The author takes responsibility for the content.
\end{document}
